\documentclass[acmsmall,screen,review,anonymous]{acmart}

\usepackage{booktabs}
\usepackage{adjustbox}
\usepackage{xcolor}
\usepackage{xspace}
\usepackage{pgfplots}
\pgfplotsset{compat=1.18}
\usepackage[most]{tcolorbox}

\newcommand{\NBugs}{580}
\newcommand{\NProjects}{74}
\newcommand{\NAgent}{198}
\newcommand{\NCandidates}{9,270}
\newcommand{\NUsable}{4,072}
\newcommand{\GoldHash}{f8b6194f}
\newcommand{\NLonger}{427}
\newcommand{\AFgpuFullAll}{77.9}
\newcommand{\AFbrowserFullAll}{75.3}
\newcommand{\AFgpuFullAgents}{76.8}
\newcommand{\LunaAgentFull}{77.3}
\newcommand{\AFgpuShortAgents}{74.2}
\newcommand{\ClaudeAgentShort}{90.9}
\newcommand{\AstraAgentShort}{86.4}
\newcommand{\AstraSingleShort}{56.1}
\newcommand{\AFgpuHitOneShort}{45.5}
\newcommand{\AstraAgentHitOneShort}{80.3}
\newcommand{\CascadeClaude}{91.4}
\newcommand{\CascadeClaudeEscalated}{62.6}
\newcommand{\CascadeClaudeRandom}{84.7}
\newcommand{\RepeatSingleChanged}{4}
\newcommand{\RepeatSingleN}{50}
\newcommand{\RepeatAgentChanged}{0}
\newcommand{\RepeatAgentN}{42}
\newcommand{\CostClaudeAgent}{\$0.20}
\newcommand{\CostAstraAgent}{\$0.29}
\newcommand{\CostAstraSingle}{\$0.51}
\newcommand{\CostAFgpu}{\$0.001}

\newcommand{\res}[3]{\ifcsname r@#1@#2@#3\endcsname\csname r@#1@#2@#3\endcsname\else\textbf{\textcolor{red}{pending}}\fi}
\newcommand{\setres}[4]{\expandafter\def\csname r@#1@#2@#3\endcsname{#4}}
\setres{fresh-full-agents}{claude-single-sonnet}{hit@1}{75.0}
\setres{fresh-full-agents}{claude-single-sonnet}{hit@1-ci}{43.8--94.1}
\setres{fresh-full-agents}{claude-single-sonnet}{hit@5}{85.0}
\setres{fresh-full-agents}{claude-single-sonnet}{hit@5-ci}{53.3--100.0}
\setres{fresh-full-agents}{claude-single-sonnet}{hit@10}{85.0}
\setres{fresh-full-agents}{claude-single-sonnet}{hit@10-ci}{53.3--100.0}
\setres{fresh-full-agents}{claude-single-sonnet}{all@10}{45.0}
\setres{fresh-full-agents}{claude-single-sonnet}{all@10-ci}{23.8--70.6}
\setres{fresh-full-agents}{claude-single-sonnet}{mrr}{0.792}
\setres{fresh-full-agents}{claude-single-sonnet}{n}{20}
\setres{fresh-full-agents}{claude-single-sonnet}{cost}{\$0.49}
\setres{fresh-full-agents}{claude-single-sonnet}{sec}{6.5}
\setres{fresh-full-agents}{codex-agent-astra}{hit@1}{79.3}
\setres{fresh-full-agents}{codex-agent-astra}{hit@1-ci}{73.2--84.9}
\setres{fresh-full-agents}{codex-agent-astra}{hit@5}{86.9}
\setres{fresh-full-agents}{codex-agent-astra}{hit@5-ci}{81.7--91.5}
\setres{fresh-full-agents}{codex-agent-astra}{hit@10}{86.9}
\setres{fresh-full-agents}{codex-agent-astra}{hit@10-ci}{81.7--91.5}
\setres{fresh-full-agents}{codex-agent-astra}{all@10}{42.4}
\setres{fresh-full-agents}{codex-agent-astra}{all@10-ci}{35.5--49.5}
\setres{fresh-full-agents}{codex-agent-astra}{mrr}{0.825}
\setres{fresh-full-agents}{codex-agent-astra}{n}{198}
\setres{fresh-full-agents}{codex-agent-astra}{cost}{\$0.28}
\setres{fresh-full-agents}{codex-agent-astra}{sec}{23.1}
\setres{fresh-full-agents}{codex-agent}{hit@1}{65.2}
\setres{fresh-full-agents}{codex-agent}{hit@1-ci}{57.1--72.6}
\setres{fresh-full-agents}{codex-agent}{hit@5}{77.3}
\setres{fresh-full-agents}{codex-agent}{hit@5-ci}{70.9--83.3}
\setres{fresh-full-agents}{codex-agent}{hit@10}{77.3}
\setres{fresh-full-agents}{codex-agent}{hit@10-ci}{70.9--83.3}
\setres{fresh-full-agents}{codex-agent}{all@10}{36.4}
\setres{fresh-full-agents}{codex-agent}{all@10-ci}{29.1--44.0}
\setres{fresh-full-agents}{codex-agent}{mrr}{0.702}
\setres{fresh-full-agents}{codex-agent}{n}{198}
\setres{fresh-full-agents}{codex-agent}{cost}{\$0.003}
\setres{fresh-full-agents}{codex-agent}{sec}{21.6}
\setres{fresh-full-agents}{codex-single-astra}{hit@1}{49.0}
\setres{fresh-full-agents}{codex-single-astra}{hit@1-ci}{40.7--57.2}
\setres{fresh-full-agents}{codex-single-astra}{hit@5}{60.1}
\setres{fresh-full-agents}{codex-single-astra}{hit@5-ci}{51.4--68.6}
\setres{fresh-full-agents}{codex-single-astra}{hit@10}{60.6}
\setres{fresh-full-agents}{codex-single-astra}{hit@10-ci}{51.9--69.0}
\setres{fresh-full-agents}{codex-single-astra}{all@10}{26.3}
\setres{fresh-full-agents}{codex-single-astra}{all@10-ci}{19.9--33.7}
\setres{fresh-full-agents}{codex-single-astra}{mrr}{0.537}
\setres{fresh-full-agents}{codex-single-astra}{n}{198}
\setres{fresh-full-agents}{codex-single-astra}{cost}{\$0.54}
\setres{fresh-full-agents}{codex-single-astra}{sec}{16.4}
\setres{fresh-full-agents}{codex-single}{hit@1}{37.4}
\setres{fresh-full-agents}{codex-single}{hit@1-ci}{28.2--47.2}
\setres{fresh-full-agents}{codex-single}{hit@5}{42.9}
\setres{fresh-full-agents}{codex-single}{hit@5-ci}{33.7--52.4}
\setres{fresh-full-agents}{codex-single}{hit@10}{42.9}
\setres{fresh-full-agents}{codex-single}{hit@10-ci}{33.7--52.4}
\setres{fresh-full-agents}{codex-single}{all@10}{20.2}
\setres{fresh-full-agents}{codex-single}{all@10-ci}{14.0--27.3}
\setres{fresh-full-agents}{codex-single}{mrr}{0.397}
\setres{fresh-full-agents}{codex-single}{n}{198}
\setres{fresh-full-agents}{codex-single}{cost}{\$0.005}
\setres{fresh-full-agents}{codex-single}{sec}{10.6}
\setres{fresh-full-agents}{gemini-single}{hit@1}{42.4}
\setres{fresh-full-agents}{gemini-single}{hit@1-ci}{33.5--51.5}
\setres{fresh-full-agents}{gemini-single}{hit@5}{53.0}
\setres{fresh-full-agents}{gemini-single}{hit@5-ci}{44.2--61.9}
\setres{fresh-full-agents}{gemini-single}{hit@10}{55.6}
\setres{fresh-full-agents}{gemini-single}{hit@10-ci}{46.3--64.9}
\setres{fresh-full-agents}{gemini-single}{all@10}{28.3}
\setres{fresh-full-agents}{gemini-single}{all@10-ci}{21.4--35.9}
\setres{fresh-full-agents}{gemini-single}{mrr}{0.468}
\setres{fresh-full-agents}{gemini-single}{n}{198}
\setres{fresh-full-agents}{gemini-single}{cost}{\$0.100}
\setres{fresh-full-agents}{gemini-single}{sec}{25.0}
\setres{fresh-full-agents}{open-agent-gpt-oss-120b}{hit@1}{48.5}
\setres{fresh-full-agents}{open-agent-gpt-oss-120b}{hit@1-ci}{39.6--56.9}
\setres{fresh-full-agents}{open-agent-gpt-oss-120b}{hit@5}{51.0}
\setres{fresh-full-agents}{open-agent-gpt-oss-120b}{hit@5-ci}{42.3--59.5}
\setres{fresh-full-agents}{open-agent-gpt-oss-120b}{hit@10}{51.0}
\setres{fresh-full-agents}{open-agent-gpt-oss-120b}{hit@10-ci}{42.3--59.5}
\setres{fresh-full-agents}{open-agent-gpt-oss-120b}{all@10}{23.7}
\setres{fresh-full-agents}{open-agent-gpt-oss-120b}{all@10-ci}{17.3--30.5}
\setres{fresh-full-agents}{open-agent-gpt-oss-120b}{mrr}{0.497}
\setres{fresh-full-agents}{open-agent-gpt-oss-120b}{n}{198}
\setres{fresh-full-agents}{open-agent-gpt-oss-120b}{cost}{\$0.049}
\setres{fresh-full-agents}{open-agent-gpt-oss-120b}{sec}{35.5}
\setres{fresh-full-agents}{open-agent-qwen3-coder-next-fp8}{hit@1}{53.5}
\setres{fresh-full-agents}{open-agent-qwen3-coder-next-fp8}{hit@1-ci}{45.0--62.4}
\setres{fresh-full-agents}{open-agent-qwen3-coder-next-fp8}{hit@5}{61.6}
\setres{fresh-full-agents}{open-agent-qwen3-coder-next-fp8}{hit@5-ci}{53.5--69.5}
\setres{fresh-full-agents}{open-agent-qwen3-coder-next-fp8}{hit@10}{61.6}
\setres{fresh-full-agents}{open-agent-qwen3-coder-next-fp8}{hit@10-ci}{53.5--69.5}
\setres{fresh-full-agents}{open-agent-qwen3-coder-next-fp8}{all@10}{34.8}
\setres{fresh-full-agents}{open-agent-qwen3-coder-next-fp8}{all@10-ci}{28.2--41.7}
\setres{fresh-full-agents}{open-agent-qwen3-coder-next-fp8}{mrr}{0.566}
\setres{fresh-full-agents}{open-agent-qwen3-coder-next-fp8}{n}{198}
\setres{fresh-full-agents}{open-agent-qwen3-coder-next-fp8}{cost}{\$0.027}
\setres{fresh-full-agents}{open-agent-qwen3-coder-next-fp8}{sec}{18.6}
\setres{fresh-full-agents}{open-single-gpt-oss-120b}{hit@1}{25.8}
\setres{fresh-full-agents}{open-single-gpt-oss-120b}{hit@1-ci}{18.0--34.1}
\setres{fresh-full-agents}{open-single-gpt-oss-120b}{hit@5}{30.3}
\setres{fresh-full-agents}{open-single-gpt-oss-120b}{hit@5-ci}{21.8--39.1}
\setres{fresh-full-agents}{open-single-gpt-oss-120b}{hit@10}{31.3}
\setres{fresh-full-agents}{open-single-gpt-oss-120b}{hit@10-ci}{22.8--40.0}
\setres{fresh-full-agents}{open-single-gpt-oss-120b}{all@10}{14.1}
\setres{fresh-full-agents}{open-single-gpt-oss-120b}{all@10-ci}{9.4--19.5}
\setres{fresh-full-agents}{open-single-gpt-oss-120b}{mrr}{0.277}
\setres{fresh-full-agents}{open-single-gpt-oss-120b}{n}{198}
\setres{fresh-full-agents}{open-single-gpt-oss-120b}{cost}{\$0.006}
\setres{fresh-full-agents}{open-single-gpt-oss-120b}{sec}{4.5}
\setres{fresh-full-agents}{open-single-qwen3-coder-next-fp8}{hit@1}{31.8}
\setres{fresh-full-agents}{open-single-qwen3-coder-next-fp8}{hit@1-ci}{23.2--41.0}
\setres{fresh-full-agents}{open-single-qwen3-coder-next-fp8}{hit@5}{38.9}
\setres{fresh-full-agents}{open-single-qwen3-coder-next-fp8}{hit@5-ci}{30.1--48.0}
\setres{fresh-full-agents}{open-single-qwen3-coder-next-fp8}{hit@10}{40.9}
\setres{fresh-full-agents}{open-single-qwen3-coder-next-fp8}{hit@10-ci}{31.7--50.0}
\setres{fresh-full-agents}{open-single-qwen3-coder-next-fp8}{all@10}{20.7}
\setres{fresh-full-agents}{open-single-qwen3-coder-next-fp8}{all@10-ci}{14.0--28.1}
\setres{fresh-full-agents}{open-single-qwen3-coder-next-fp8}{mrr}{0.351}
\setres{fresh-full-agents}{open-single-qwen3-coder-next-fp8}{n}{198}
\setres{fresh-full-agents}{open-single-qwen3-coder-next-fp8}{cost}{\$0.006}
\setres{fresh-full-agents}{open-single-qwen3-coder-next-fp8}{sec}{3.0}
\setres{fresh-full-agents}{swerank-embed-large}{hit@1}{40.4}
\setres{fresh-full-agents}{swerank-embed-large}{hit@1-ci}{33.3--47.5}
\setres{fresh-full-agents}{swerank-embed-large}{hit@5}{65.2}
\setres{fresh-full-agents}{swerank-embed-large}{hit@5-ci}{58.3--72.2}
\setres{fresh-full-agents}{swerank-embed-large}{hit@10}{72.7}
\setres{fresh-full-agents}{swerank-embed-large}{hit@10-ci}{66.1--79.6}
\setres{fresh-full-agents}{swerank-embed-large}{all@10}{38.4}
\setres{fresh-full-agents}{swerank-embed-large}{all@10-ci}{30.4--47.1}
\setres{fresh-full-agents}{swerank-embed-large}{mrr}{0.518}
\setres{fresh-full-agents}{swerank-embed-large}{n}{198}
\setres{fresh-full-agents}{swerank-embed-large}{cost}{--}
\setres{fresh-full-agents}{swerank-embed-large}{sec}{0.0}
\setres{fresh-full-agents}{swerank-embed-small}{hit@1}{38.4}
\setres{fresh-full-agents}{swerank-embed-small}{hit@1-ci}{30.9--46.2}
\setres{fresh-full-agents}{swerank-embed-small}{hit@5}{57.1}
\setres{fresh-full-agents}{swerank-embed-small}{hit@5-ci}{49.5--64.4}
\setres{fresh-full-agents}{swerank-embed-small}{hit@10}{62.1}
\setres{fresh-full-agents}{swerank-embed-small}{hit@10-ci}{54.5--69.7}
\setres{fresh-full-agents}{swerank-embed-small}{all@10}{33.8}
\setres{fresh-full-agents}{swerank-embed-small}{all@10-ci}{27.7--40.5}
\setres{fresh-full-agents}{swerank-embed-small}{mrr}{0.459}
\setres{fresh-full-agents}{swerank-embed-small}{n}{198}
\setres{fresh-full-agents}{swerank-embed-small}{cost}{--}
\setres{fresh-full-agents}{swerank-embed-small}{sec}{0.0}
\setres{fresh-full-agents}{website-gpu}{hit@1}{47.5}
\setres{fresh-full-agents}{website-gpu}{hit@1-ci}{40.5--54.5}
\setres{fresh-full-agents}{website-gpu}{hit@5}{69.7}
\setres{fresh-full-agents}{website-gpu}{hit@5-ci}{62.6--76.5}
\setres{fresh-full-agents}{website-gpu}{hit@10}{76.8}
\setres{fresh-full-agents}{website-gpu}{hit@10-ci}{69.8--83.0}
\setres{fresh-full-agents}{website-gpu}{all@10}{45.5}
\setres{fresh-full-agents}{website-gpu}{all@10-ci}{37.6--53.1}
\setres{fresh-full-agents}{website-gpu}{mrr}{0.566}
\setres{fresh-full-agents}{website-gpu}{n}{198}
\setres{fresh-full-agents}{website-gpu}{cost}{\$0.001}
\setres{fresh-full-agents}{website-gpu}{sec}{8.4}
\setres{fresh-full-agents}{website}{hit@1}{46.0}
\setres{fresh-full-agents}{website}{hit@1-ci}{39.1--52.8}
\setres{fresh-full-agents}{website}{hit@5}{69.2}
\setres{fresh-full-agents}{website}{hit@5-ci}{62.6--75.5}
\setres{fresh-full-agents}{website}{hit@10}{77.3}
\setres{fresh-full-agents}{website}{hit@10-ci}{71.0--83.1}
\setres{fresh-full-agents}{website}{all@10}{42.9}
\setres{fresh-full-agents}{website}{all@10-ci}{35.5--50.7}
\setres{fresh-full-agents}{website}{mrr}{0.566}
\setres{fresh-full-agents}{website}{n}{198}
\setres{fresh-full-agents}{website}{cost}{\$0}
\setres{fresh-full-agents}{website}{sec}{3.0}
\setres{fresh-full-agents}{claude-single-sonnet}{dhit@1}{+15.0}
\setres{fresh-full-agents}{claude-single-sonnet}{wlhit@1}{4/1}
\setres{fresh-full-agents}{claude-single-sonnet}{phit@1}{=1.000}
\setres{fresh-full-agents}{claude-single-sonnet}{dhit@10}{+5.0}
\setres{fresh-full-agents}{claude-single-sonnet}{wlhit@10}{2/1}
\setres{fresh-full-agents}{claude-single-sonnet}{phit@10}{=1.000}
\setres{fresh-full-agents}{codex-agent-astra}{dhit@1}{+31.8}
\setres{fresh-full-agents}{codex-agent-astra}{wlhit@1}{73/10}
\setres{fresh-full-agents}{codex-agent-astra}{phit@1}{<0.001}
\setres{fresh-full-agents}{codex-agent-astra}{dhit@10}{+10.1}
\setres{fresh-full-agents}{codex-agent-astra}{wlhit@10}{35/15}
\setres{fresh-full-agents}{codex-agent-astra}{phit@10}{=0.198}
\setres{fresh-full-agents}{codex-agent}{dhit@1}{+17.7}
\setres{fresh-full-agents}{codex-agent}{wlhit@1}{53/18}
\setres{fresh-full-agents}{codex-agent}{phit@1}{<0.001}
\setres{fresh-full-agents}{codex-agent}{dhit@10}{+0.5}
\setres{fresh-full-agents}{codex-agent}{wlhit@10}{26/25}
\setres{fresh-full-agents}{codex-agent}{phit@10}{=1.000}
\setres{fresh-full-agents}{codex-single-astra}{dhit@1}{+1.5}
\setres{fresh-full-agents}{codex-single-astra}{wlhit@1}{35/32}
\setres{fresh-full-agents}{codex-single-astra}{phit@1}{=1.000}
\setres{fresh-full-agents}{codex-single-astra}{dhit@10}{-16.2}
\setres{fresh-full-agents}{codex-single-astra}{wlhit@10}{16/48}
\setres{fresh-full-agents}{codex-single-astra}{phit@10}{=0.005}
\setres{fresh-full-agents}{codex-single}{dhit@1}{-10.1}
\setres{fresh-full-agents}{codex-single}{wlhit@1}{17/37}
\setres{fresh-full-agents}{codex-single}{phit@1}{=0.264}
\setres{fresh-full-agents}{codex-single}{dhit@10}{-33.8}
\setres{fresh-full-agents}{codex-single}{wlhit@10}{4/71}
\setres{fresh-full-agents}{codex-single}{phit@10}{<0.001}
\setres{fresh-full-agents}{gemini-single}{dhit@1}{-5.1}
\setres{fresh-full-agents}{gemini-single}{wlhit@1}{25/35}
\setres{fresh-full-agents}{gemini-single}{phit@1}{=1.000}
\setres{fresh-full-agents}{gemini-single}{dhit@10}{-21.2}
\setres{fresh-full-agents}{gemini-single}{wlhit@10}{11/53}
\setres{fresh-full-agents}{gemini-single}{phit@10}{<0.001}
\setres{fresh-full-agents}{open-agent-gpt-oss-120b}{dhit@1}{+1.0}
\setres{fresh-full-agents}{open-agent-gpt-oss-120b}{wlhit@1}{31/29}
\setres{fresh-full-agents}{open-agent-gpt-oss-120b}{phit@1}{=1.000}
\setres{fresh-full-agents}{open-agent-gpt-oss-120b}{dhit@10}{-25.8}
\setres{fresh-full-agents}{open-agent-gpt-oss-120b}{wlhit@10}{10/61}
\setres{fresh-full-agents}{open-agent-gpt-oss-120b}{phit@10}{<0.001}
\setres{fresh-full-agents}{open-agent-qwen3-coder-next-fp8}{dhit@1}{+6.1}
\setres{fresh-full-agents}{open-agent-qwen3-coder-next-fp8}{wlhit@1}{38/26}
\setres{fresh-full-agents}{open-agent-qwen3-coder-next-fp8}{phit@1}{=1.000}
\setres{fresh-full-agents}{open-agent-qwen3-coder-next-fp8}{dhit@10}{-15.2}
\setres{fresh-full-agents}{open-agent-qwen3-coder-next-fp8}{wlhit@10}{15/45}
\setres{fresh-full-agents}{open-agent-qwen3-coder-next-fp8}{phit@10}{=0.005}
\setres{fresh-full-agents}{open-single-gpt-oss-120b}{dhit@1}{-21.7}
\setres{fresh-full-agents}{open-single-gpt-oss-120b}{wlhit@1}{9/52}
\setres{fresh-full-agents}{open-single-gpt-oss-120b}{phit@1}{<0.001}
\setres{fresh-full-agents}{open-single-gpt-oss-120b}{dhit@10}{-45.5}
\setres{fresh-full-agents}{open-single-gpt-oss-120b}{wlhit@10}{3/93}
\setres{fresh-full-agents}{open-single-gpt-oss-120b}{phit@10}{<0.001}
\setres{fresh-full-agents}{open-single-qwen3-coder-next-fp8}{dhit@1}{-15.7}
\setres{fresh-full-agents}{open-single-qwen3-coder-next-fp8}{wlhit@1}{18/49}
\setres{fresh-full-agents}{open-single-qwen3-coder-next-fp8}{phit@1}{=0.007}
\setres{fresh-full-agents}{open-single-qwen3-coder-next-fp8}{dhit@10}{-35.9}
\setres{fresh-full-agents}{open-single-qwen3-coder-next-fp8}{wlhit@10}{5/76}
\setres{fresh-full-agents}{open-single-qwen3-coder-next-fp8}{phit@10}{<0.001}
\setres{fresh-full-agents}{swerank-embed-large}{dhit@1}{-7.1}
\setres{fresh-full-agents}{swerank-embed-large}{wlhit@1}{26/40}
\setres{fresh-full-agents}{swerank-embed-large}{phit@1}{=1.000}
\setres{fresh-full-agents}{swerank-embed-large}{dhit@10}{-4.0}
\setres{fresh-full-agents}{swerank-embed-large}{wlhit@10}{14/22}
\setres{fresh-full-agents}{swerank-embed-large}{phit@10}{=1.000}
\setres{fresh-full-agents}{swerank-embed-small}{dhit@1}{-9.1}
\setres{fresh-full-agents}{swerank-embed-small}{wlhit@1}{18/36}
\setres{fresh-full-agents}{swerank-embed-small}{phit@1}{=0.535}
\setres{fresh-full-agents}{swerank-embed-small}{dhit@10}{-14.6}
\setres{fresh-full-agents}{swerank-embed-small}{wlhit@10}{11/40}
\setres{fresh-full-agents}{swerank-embed-small}{phit@10}{=0.005}
\setres{fresh-full-agents}{website}{dhit@1}{-1.5}
\setres{fresh-full-agents}{website}{wlhit@1}{26/29}
\setres{fresh-full-agents}{website}{phit@1}{=1.000}
\setres{fresh-full-agents}{website}{dhit@10}{+0.5}
\setres{fresh-full-agents}{website}{wlhit@10}{7/6}
\setres{fresh-full-agents}{website}{phit@10}{=1.000}
\setres{fresh-full-agents}{claude-single-sonnet}{browser-dhit@1}{+10.0}
\setres{fresh-full-agents}{claude-single-sonnet}{browser-wlhit@1}{3/1}
\setres{fresh-full-agents}{claude-single-sonnet}{browser-phit@1}{=1.000}
\setres{fresh-full-agents}{claude-single-sonnet}{browser-dhit@10}{+5.0}
\setres{fresh-full-agents}{claude-single-sonnet}{browser-wlhit@10}{2/1}
\setres{fresh-full-agents}{claude-single-sonnet}{browser-phit@10}{=1.000}
\setres{fresh-full-agents}{codex-agent-astra}{browser-dhit@1}{+33.3}
\setres{fresh-full-agents}{codex-agent-astra}{browser-wlhit@1}{78/12}
\setres{fresh-full-agents}{codex-agent-astra}{browser-phit@1}{<0.001}
\setres{fresh-full-agents}{codex-agent-astra}{browser-dhit@10}{+9.6}
\setres{fresh-full-agents}{codex-agent-astra}{browser-wlhit@10}{33/14}
\setres{fresh-full-agents}{codex-agent-astra}{browser-phit@10}{=0.253}
\setres{fresh-full-agents}{codex-agent}{browser-dhit@1}{+19.2}
\setres{fresh-full-agents}{codex-agent}{browser-wlhit@1}{59/21}
\setres{fresh-full-agents}{codex-agent}{browser-phit@1}{<0.001}
\setres{fresh-full-agents}{codex-agent}{browser-dhit@10}{+0.0}
\setres{fresh-full-agents}{codex-agent}{browser-wlhit@10}{26/26}
\setres{fresh-full-agents}{codex-agent}{browser-phit@10}{=1.000}
\setres{fresh-full-agents}{codex-single-astra}{browser-dhit@1}{+3.0}
\setres{fresh-full-agents}{codex-single-astra}{browser-wlhit@1}{38/32}
\setres{fresh-full-agents}{codex-single-astra}{browser-phit@1}{=1.000}
\setres{fresh-full-agents}{codex-single-astra}{browser-dhit@10}{-16.7}
\setres{fresh-full-agents}{codex-single-astra}{browser-wlhit@10}{15/48}
\setres{fresh-full-agents}{codex-single-astra}{browser-phit@10}{<0.001}
\setres{fresh-full-agents}{codex-single}{browser-dhit@1}{-8.6}
\setres{fresh-full-agents}{codex-single}{browser-wlhit@1}{24/41}
\setres{fresh-full-agents}{codex-single}{browser-phit@1}{=1.000}
\setres{fresh-full-agents}{codex-single}{browser-dhit@10}{-34.3}
\setres{fresh-full-agents}{codex-single}{browser-wlhit@10}{4/72}
\setres{fresh-full-agents}{codex-single}{browser-phit@10}{<0.001}
\setres{fresh-full-agents}{gemini-single}{browser-dhit@1}{-3.5}
\setres{fresh-full-agents}{gemini-single}{browser-wlhit@1}{27/34}
\setres{fresh-full-agents}{gemini-single}{browser-phit@1}{=1.000}
\setres{fresh-full-agents}{gemini-single}{browser-dhit@10}{-21.7}
\setres{fresh-full-agents}{gemini-single}{browser-wlhit@10}{9/52}
\setres{fresh-full-agents}{gemini-single}{browser-phit@10}{<0.001}
\setres{fresh-full-agents}{open-agent-gpt-oss-120b}{browser-dhit@1}{+2.5}
\setres{fresh-full-agents}{open-agent-gpt-oss-120b}{browser-wlhit@1}{37/32}
\setres{fresh-full-agents}{open-agent-gpt-oss-120b}{browser-phit@1}{=1.000}
\setres{fresh-full-agents}{open-agent-gpt-oss-120b}{browser-dhit@10}{-26.3}
\setres{fresh-full-agents}{open-agent-gpt-oss-120b}{browser-wlhit@10}{7/59}
\setres{fresh-full-agents}{open-agent-gpt-oss-120b}{browser-phit@10}{<0.001}
\setres{fresh-full-agents}{open-agent-qwen3-coder-next-fp8}{browser-dhit@1}{+7.6}
\setres{fresh-full-agents}{open-agent-qwen3-coder-next-fp8}{browser-wlhit@1}{44/29}
\setres{fresh-full-agents}{open-agent-qwen3-coder-next-fp8}{browser-phit@1}{=1.000}
\setres{fresh-full-agents}{open-agent-qwen3-coder-next-fp8}{browser-dhit@10}{-15.7}
\setres{fresh-full-agents}{open-agent-qwen3-coder-next-fp8}{browser-wlhit@10}{13/44}
\setres{fresh-full-agents}{open-agent-qwen3-coder-next-fp8}{browser-phit@10}{<0.001}
\setres{fresh-full-agents}{open-single-gpt-oss-120b}{browser-dhit@1}{-20.2}
\setres{fresh-full-agents}{open-single-gpt-oss-120b}{browser-wlhit@1}{13/53}
\setres{fresh-full-agents}{open-single-gpt-oss-120b}{browser-phit@1}{<0.001}
\setres{fresh-full-agents}{open-single-gpt-oss-120b}{browser-dhit@10}{-46.0}
\setres{fresh-full-agents}{open-single-gpt-oss-120b}{browser-wlhit@10}{2/93}
\setres{fresh-full-agents}{open-single-gpt-oss-120b}{browser-phit@10}{<0.001}
\setres{fresh-full-agents}{open-single-qwen3-coder-next-fp8}{browser-dhit@1}{-14.1}
\setres{fresh-full-agents}{open-single-qwen3-coder-next-fp8}{browser-wlhit@1}{20/48}
\setres{fresh-full-agents}{open-single-qwen3-coder-next-fp8}{browser-phit@1}{=0.033}
\setres{fresh-full-agents}{open-single-qwen3-coder-next-fp8}{browser-dhit@10}{-36.4}
\setres{fresh-full-agents}{open-single-qwen3-coder-next-fp8}{browser-wlhit@10}{4/76}
\setres{fresh-full-agents}{open-single-qwen3-coder-next-fp8}{browser-phit@10}{<0.001}
\setres{fresh-full-agents}{swerank-embed-large}{browser-dhit@1}{-5.6}
\setres{fresh-full-agents}{swerank-embed-large}{browser-wlhit@1}{27/38}
\setres{fresh-full-agents}{swerank-embed-large}{browser-phit@1}{=1.000}
\setres{fresh-full-agents}{swerank-embed-large}{browser-dhit@10}{-4.5}
\setres{fresh-full-agents}{swerank-embed-large}{browser-wlhit@10}{15/24}
\setres{fresh-full-agents}{swerank-embed-large}{browser-phit@10}{=1.000}
\setres{fresh-full-agents}{swerank-embed-small}{browser-dhit@1}{-7.6}
\setres{fresh-full-agents}{swerank-embed-small}{browser-wlhit@1}{24/39}
\setres{fresh-full-agents}{swerank-embed-small}{browser-phit@1}{=1.000}
\setres{fresh-full-agents}{swerank-embed-small}{browser-dhit@10}{-15.2}
\setres{fresh-full-agents}{swerank-embed-small}{browser-wlhit@10}{11/41}
\setres{fresh-full-agents}{swerank-embed-small}{browser-phit@10}{<0.001}
\setres{fresh-full-agents}{website-gpu}{browser-dhit@1}{+1.5}
\setres{fresh-full-agents}{website-gpu}{browser-wlhit@1}{29/26}
\setres{fresh-full-agents}{website-gpu}{browser-phit@1}{=1.000}
\setres{fresh-full-agents}{website-gpu}{browser-dhit@10}{-0.5}
\setres{fresh-full-agents}{website-gpu}{browser-wlhit@10}{6/7}
\setres{fresh-full-agents}{website-gpu}{browser-phit@10}{=1.000}
\setres{fresh-full-agents}{claude-single-sonnet}{s-function}{86.7}
\setres{fresh-full-agents}{claude-single-sonnet}{sn-function}{15}
\setres{fresh-full-agents}{claude-single-sonnet}{s-file}{100.0}
\setres{fresh-full-agents}{claude-single-sonnet}{sn-file}{1}
\setres{fresh-full-agents}{claude-single-sonnet}{s-none}{75.0}
\setres{fresh-full-agents}{claude-single-sonnet}{sn-none}{4}
\setres{fresh-full-agents}{codex-agent-astra}{s-function}{91.9}
\setres{fresh-full-agents}{codex-agent-astra}{sn-function}{86}
\setres{fresh-full-agents}{codex-agent-astra}{s-file}{87.0}
\setres{fresh-full-agents}{codex-agent-astra}{sn-file}{23}
\setres{fresh-full-agents}{codex-agent-astra}{s-trace}{75.0}
\setres{fresh-full-agents}{codex-agent-astra}{sn-trace}{4}
\setres{fresh-full-agents}{codex-agent-astra}{s-none}{82.3}
\setres{fresh-full-agents}{codex-agent-astra}{sn-none}{85}
\setres{fresh-full-agents}{codex-agent}{s-function}{87.2}
\setres{fresh-full-agents}{codex-agent}{sn-function}{86}
\setres{fresh-full-agents}{codex-agent}{s-file}{82.6}
\setres{fresh-full-agents}{codex-agent}{sn-file}{23}
\setres{fresh-full-agents}{codex-agent}{s-trace}{50.0}
\setres{fresh-full-agents}{codex-agent}{sn-trace}{4}
\setres{fresh-full-agents}{codex-agent}{s-none}{67.1}
\setres{fresh-full-agents}{codex-agent}{sn-none}{85}
\setres{fresh-full-agents}{codex-single-astra}{s-function}{87.2}
\setres{fresh-full-agents}{codex-single-astra}{sn-function}{86}
\setres{fresh-full-agents}{codex-single-astra}{s-file}{69.6}
\setres{fresh-full-agents}{codex-single-astra}{sn-file}{23}
\setres{fresh-full-agents}{codex-single-astra}{s-trace}{0.0}
\setres{fresh-full-agents}{codex-single-astra}{sn-trace}{4}
\setres{fresh-full-agents}{codex-single-astra}{s-none}{34.1}
\setres{fresh-full-agents}{codex-single-astra}{sn-none}{85}
\setres{fresh-full-agents}{codex-single}{s-function}{75.6}
\setres{fresh-full-agents}{codex-single}{sn-function}{86}
\setres{fresh-full-agents}{codex-single}{s-file}{43.5}
\setres{fresh-full-agents}{codex-single}{sn-file}{23}
\setres{fresh-full-agents}{codex-single}{s-trace}{0.0}
\setres{fresh-full-agents}{codex-single}{sn-trace}{4}
\setres{fresh-full-agents}{codex-single}{s-none}{11.8}
\setres{fresh-full-agents}{codex-single}{sn-none}{85}
\setres{fresh-full-agents}{gemini-single}{s-function}{86.1}
\setres{fresh-full-agents}{gemini-single}{sn-function}{86}
\setres{fresh-full-agents}{gemini-single}{s-file}{52.2}
\setres{fresh-full-agents}{gemini-single}{sn-file}{23}
\setres{fresh-full-agents}{gemini-single}{s-trace}{0.0}
\setres{fresh-full-agents}{gemini-single}{sn-trace}{4}
\setres{fresh-full-agents}{gemini-single}{s-none}{28.2}
\setres{fresh-full-agents}{gemini-single}{sn-none}{85}
\setres{fresh-full-agents}{open-agent-gpt-oss-120b}{s-function}{72.1}
\setres{fresh-full-agents}{open-agent-gpt-oss-120b}{sn-function}{86}
\setres{fresh-full-agents}{open-agent-gpt-oss-120b}{s-file}{60.9}
\setres{fresh-full-agents}{open-agent-gpt-oss-120b}{sn-file}{23}
\setres{fresh-full-agents}{open-agent-gpt-oss-120b}{s-trace}{25.0}
\setres{fresh-full-agents}{open-agent-gpt-oss-120b}{sn-trace}{4}
\setres{fresh-full-agents}{open-agent-gpt-oss-120b}{s-none}{28.2}
\setres{fresh-full-agents}{open-agent-gpt-oss-120b}{sn-none}{85}
\setres{fresh-full-agents}{open-agent-qwen3-coder-next-fp8}{s-function}{84.9}
\setres{fresh-full-agents}{open-agent-qwen3-coder-next-fp8}{sn-function}{86}
\setres{fresh-full-agents}{open-agent-qwen3-coder-next-fp8}{s-file}{65.2}
\setres{fresh-full-agents}{open-agent-qwen3-coder-next-fp8}{sn-file}{23}
\setres{fresh-full-agents}{open-agent-qwen3-coder-next-fp8}{s-trace}{25.0}
\setres{fresh-full-agents}{open-agent-qwen3-coder-next-fp8}{sn-trace}{4}
\setres{fresh-full-agents}{open-agent-qwen3-coder-next-fp8}{s-none}{38.8}
\setres{fresh-full-agents}{open-agent-qwen3-coder-next-fp8}{sn-none}{85}
\setres{fresh-full-agents}{open-single-gpt-oss-120b}{s-function}{59.3}
\setres{fresh-full-agents}{open-single-gpt-oss-120b}{sn-function}{86}
\setres{fresh-full-agents}{open-single-gpt-oss-120b}{s-file}{26.1}
\setres{fresh-full-agents}{open-single-gpt-oss-120b}{sn-file}{23}
\setres{fresh-full-agents}{open-single-gpt-oss-120b}{s-trace}{0.0}
\setres{fresh-full-agents}{open-single-gpt-oss-120b}{sn-trace}{4}
\setres{fresh-full-agents}{open-single-gpt-oss-120b}{s-none}{5.9}
\setres{fresh-full-agents}{open-single-gpt-oss-120b}{sn-none}{85}
\setres{fresh-full-agents}{open-single-qwen3-coder-next-fp8}{s-function}{73.3}
\setres{fresh-full-agents}{open-single-qwen3-coder-next-fp8}{sn-function}{86}
\setres{fresh-full-agents}{open-single-qwen3-coder-next-fp8}{s-file}{30.4}
\setres{fresh-full-agents}{open-single-qwen3-coder-next-fp8}{sn-file}{23}
\setres{fresh-full-agents}{open-single-qwen3-coder-next-fp8}{s-trace}{0.0}
\setres{fresh-full-agents}{open-single-qwen3-coder-next-fp8}{sn-trace}{4}
\setres{fresh-full-agents}{open-single-qwen3-coder-next-fp8}{s-none}{12.9}
\setres{fresh-full-agents}{open-single-qwen3-coder-next-fp8}{sn-none}{85}
\setres{fresh-full-agents}{swerank-embed-large}{s-function}{84.9}
\setres{fresh-full-agents}{swerank-embed-large}{sn-function}{86}
\setres{fresh-full-agents}{swerank-embed-large}{s-file}{73.9}
\setres{fresh-full-agents}{swerank-embed-large}{sn-file}{23}
\setres{fresh-full-agents}{swerank-embed-large}{s-trace}{0.0}
\setres{fresh-full-agents}{swerank-embed-large}{sn-trace}{4}
\setres{fresh-full-agents}{swerank-embed-large}{s-none}{63.5}
\setres{fresh-full-agents}{swerank-embed-large}{sn-none}{85}
\setres{fresh-full-agents}{swerank-embed-small}{s-function}{69.8}
\setres{fresh-full-agents}{swerank-embed-small}{sn-function}{86}
\setres{fresh-full-agents}{swerank-embed-small}{s-file}{65.2}
\setres{fresh-full-agents}{swerank-embed-small}{sn-file}{23}
\setres{fresh-full-agents}{swerank-embed-small}{s-trace}{0.0}
\setres{fresh-full-agents}{swerank-embed-small}{sn-trace}{4}
\setres{fresh-full-agents}{swerank-embed-small}{s-none}{56.5}
\setres{fresh-full-agents}{swerank-embed-small}{sn-none}{85}
\setres{fresh-full-agents}{website-gpu}{s-function}{94.2}
\setres{fresh-full-agents}{website-gpu}{sn-function}{86}
\setres{fresh-full-agents}{website-gpu}{s-file}{73.9}
\setres{fresh-full-agents}{website-gpu}{sn-file}{23}
\setres{fresh-full-agents}{website-gpu}{s-trace}{0.0}
\setres{fresh-full-agents}{website-gpu}{sn-trace}{4}
\setres{fresh-full-agents}{website-gpu}{s-none}{63.5}
\setres{fresh-full-agents}{website-gpu}{sn-none}{85}
\setres{fresh-full-agents}{website}{s-function}{95.3}
\setres{fresh-full-agents}{website}{sn-function}{86}
\setres{fresh-full-agents}{website}{s-file}{78.3}
\setres{fresh-full-agents}{website}{sn-file}{23}
\setres{fresh-full-agents}{website}{s-trace}{0.0}
\setres{fresh-full-agents}{website}{sn-trace}{4}
\setres{fresh-full-agents}{website}{s-none}{62.4}
\setres{fresh-full-agents}{website}{sn-none}{85}
\setres{fresh-full-agents}{claude-single-sonnet}{esc-escalated}{25.0}
\setres{fresh-full-agents}{claude-single-sonnet}{esc-hit@10}{90.0}
\setres{fresh-full-agents}{claude-single-sonnet}{esc-random_hit@10}{81.2}
\setres{fresh-full-agents}{codex-agent-astra}{esc-escalated}{56.1}
\setres{fresh-full-agents}{codex-agent-astra}{esc-hit@10}{89.9}
\setres{fresh-full-agents}{codex-agent-astra}{esc-random_hit@10}{82.4}
\setres{fresh-full-agents}{codex-agent}{esc-escalated}{56.1}
\setres{fresh-full-agents}{codex-agent}{esc-hit@10}{81.8}
\setres{fresh-full-agents}{codex-agent}{esc-random_hit@10}{77.0}
\setres{fresh-full-agents}{codex-single-astra}{esc-escalated}{56.1}
\setres{fresh-full-agents}{codex-single-astra}{esc-hit@10}{68.2}
\setres{fresh-full-agents}{codex-single-astra}{esc-random_hit@10}{67.7}
\setres{fresh-full-agents}{codex-single}{esc-escalated}{56.1}
\setres{fresh-full-agents}{codex-single}{esc-hit@10}{54.0}
\setres{fresh-full-agents}{codex-single}{esc-random_hit@10}{57.8}
\setres{fresh-full-agents}{gemini-single}{esc-escalated}{56.1}
\setres{fresh-full-agents}{gemini-single}{esc-hit@10}{62.1}
\setres{fresh-full-agents}{gemini-single}{esc-random_hit@10}{64.9}
\setres{fresh-full-agents}{open-agent-gpt-oss-120b}{esc-escalated}{56.1}
\setres{fresh-full-agents}{open-agent-gpt-oss-120b}{esc-hit@10}{63.6}
\setres{fresh-full-agents}{open-agent-gpt-oss-120b}{esc-random_hit@10}{62.3}
\setres{fresh-full-agents}{open-agent-qwen3-coder-next-fp8}{esc-escalated}{56.1}
\setres{fresh-full-agents}{open-agent-qwen3-coder-next-fp8}{esc-hit@10}{67.2}
\setres{fresh-full-agents}{open-agent-qwen3-coder-next-fp8}{esc-random_hit@10}{68.3}
\setres{fresh-full-agents}{open-single-gpt-oss-120b}{esc-escalated}{56.1}
\setres{fresh-full-agents}{open-single-gpt-oss-120b}{esc-hit@10}{49.0}
\setres{fresh-full-agents}{open-single-gpt-oss-120b}{esc-random_hit@10}{51.3}
\setres{fresh-full-agents}{open-single-qwen3-coder-next-fp8}{esc-escalated}{56.1}
\setres{fresh-full-agents}{open-single-qwen3-coder-next-fp8}{esc-hit@10}{53.5}
\setres{fresh-full-agents}{open-single-qwen3-coder-next-fp8}{esc-random_hit@10}{56.7}
\setres{fresh-full-agents}{swerank-embed-large}{esc-escalated}{56.1}
\setres{fresh-full-agents}{swerank-embed-large}{esc-hit@10}{77.3}
\setres{fresh-full-agents}{swerank-embed-large}{esc-random_hit@10}{74.5}
\setres{fresh-full-agents}{swerank-embed-small}{esc-escalated}{56.1}
\setres{fresh-full-agents}{swerank-embed-small}{esc-hit@10}{70.2}
\setres{fresh-full-agents}{swerank-embed-small}{esc-random_hit@10}{68.6}
\setres{fresh-full-agents}{website}{esc-escalated}{56.1}
\setres{fresh-full-agents}{website}{esc-hit@10}{76.3}
\setres{fresh-full-agents}{website}{esc-random_hit@10}{77.0}
\setres{fresh-full}{codex-single-astra}{hit@1}{45.0}
\setres{fresh-full}{codex-single-astra}{hit@1-ci}{39.1--50.4}
\setres{fresh-full}{codex-single-astra}{hit@5}{54.7}
\setres{fresh-full}{codex-single-astra}{hit@5-ci}{48.3--60.5}
\setres{fresh-full}{codex-single-astra}{hit@10}{55.0}
\setres{fresh-full}{codex-single-astra}{hit@10-ci}{48.5--60.9}
\setres{fresh-full}{codex-single-astra}{all@10}{24.8}
\setres{fresh-full}{codex-single-astra}{all@10-ci}{20.0--29.7}
\setres{fresh-full}{codex-single-astra}{mrr}{0.492}
\setres{fresh-full}{codex-single-astra}{n}{580}
\setres{fresh-full}{codex-single-astra}{cost}{\$0.50}
\setres{fresh-full}{codex-single-astra}{sec}{15.6}
\setres{fresh-full}{codex-single}{hit@1}{34.0}
\setres{fresh-full}{codex-single}{hit@1-ci}{27.9--40.3}
\setres{fresh-full}{codex-single}{hit@5}{39.8}
\setres{fresh-full}{codex-single}{hit@5-ci}{33.3--46.3}
\setres{fresh-full}{codex-single}{hit@10}{39.8}
\setres{fresh-full}{codex-single}{hit@10-ci}{33.3--46.3}
\setres{fresh-full}{codex-single}{all@10}{19.1}
\setres{fresh-full}{codex-single}{all@10-ci}{14.7--23.8}
\setres{fresh-full}{codex-single}{mrr}{0.366}
\setres{fresh-full}{codex-single}{n}{580}
\setres{fresh-full}{codex-single}{cost}{\$0.005}
\setres{fresh-full}{codex-single}{sec}{10.5}
\setres{fresh-full}{gemini-single}{hit@1}{40.9}
\setres{fresh-full}{gemini-single}{hit@1-ci}{33.8--47.9}
\setres{fresh-full}{gemini-single}{hit@5}{49.4}
\setres{fresh-full}{gemini-single}{hit@5-ci}{42.4--56.3}
\setres{fresh-full}{gemini-single}{hit@10}{52.7}
\setres{fresh-full}{gemini-single}{hit@10-ci}{45.6--59.6}
\setres{fresh-full}{gemini-single}{all@10}{26.8}
\setres{fresh-full}{gemini-single}{all@10-ci}{21.1--32.7}
\setres{fresh-full}{gemini-single}{mrr}{0.446}
\setres{fresh-full}{gemini-single}{n}{425}
\setres{fresh-full}{gemini-single}{cost}{\$0.11}
\setres{fresh-full}{gemini-single}{sec}{28.3}
\setres{fresh-full}{open-single-gpt-oss-120b}{hit@1}{21.7}
\setres{fresh-full}{open-single-gpt-oss-120b}{hit@1-ci}{16.4--27.2}
\setres{fresh-full}{open-single-gpt-oss-120b}{hit@5}{26.6}
\setres{fresh-full}{open-single-gpt-oss-120b}{hit@5-ci}{20.4--32.8}
\setres{fresh-full}{open-single-gpt-oss-120b}{hit@10}{27.2}
\setres{fresh-full}{open-single-gpt-oss-120b}{hit@10-ci}{21.1--33.5}
\setres{fresh-full}{open-single-gpt-oss-120b}{all@10}{13.6}
\setres{fresh-full}{open-single-gpt-oss-120b}{all@10-ci}{9.5--18.0}
\setres{fresh-full}{open-single-gpt-oss-120b}{mrr}{0.239}
\setres{fresh-full}{open-single-gpt-oss-120b}{n}{580}
\setres{fresh-full}{open-single-gpt-oss-120b}{cost}{\$0.008}
\setres{fresh-full}{open-single-gpt-oss-120b}{sec}{5.0}
\setres{fresh-full}{open-single-qwen3-coder-next-fp8}{hit@1}{29.0}
\setres{fresh-full}{open-single-qwen3-coder-next-fp8}{hit@1-ci}{23.6--34.6}
\setres{fresh-full}{open-single-qwen3-coder-next-fp8}{hit@5}{35.2}
\setres{fresh-full}{open-single-qwen3-coder-next-fp8}{hit@5-ci}{29.2--41.1}
\setres{fresh-full}{open-single-qwen3-coder-next-fp8}{hit@10}{36.7}
\setres{fresh-full}{open-single-qwen3-coder-next-fp8}{hit@10-ci}{30.6--42.8}
\setres{fresh-full}{open-single-qwen3-coder-next-fp8}{all@10}{18.3}
\setres{fresh-full}{open-single-qwen3-coder-next-fp8}{all@10-ci}{13.9--23.0}
\setres{fresh-full}{open-single-qwen3-coder-next-fp8}{mrr}{0.318}
\setres{fresh-full}{open-single-qwen3-coder-next-fp8}{n}{580}
\setres{fresh-full}{open-single-qwen3-coder-next-fp8}{cost}{\$0.006}
\setres{fresh-full}{open-single-qwen3-coder-next-fp8}{sec}{3.0}
\setres{fresh-full}{swerank-embed-large}{hit@1}{43.3}
\setres{fresh-full}{swerank-embed-large}{hit@1-ci}{37.9--48.4}
\setres{fresh-full}{swerank-embed-large}{hit@5}{66.9}
\setres{fresh-full}{swerank-embed-large}{hit@5-ci}{60.6--72.6}
\setres{fresh-full}{swerank-embed-large}{hit@10}{72.9}
\setres{fresh-full}{swerank-embed-large}{hit@10-ci}{66.5--78.5}
\setres{fresh-full}{swerank-embed-large}{all@10}{37.9}
\setres{fresh-full}{swerank-embed-large}{all@10-ci}{32.0--43.6}
\setres{fresh-full}{swerank-embed-large}{mrr}{0.538}
\setres{fresh-full}{swerank-embed-large}{n}{580}
\setres{fresh-full}{swerank-embed-large}{cost}{--}
\setres{fresh-full}{swerank-embed-large}{sec}{0.0}
\setres{fresh-full}{swerank-embed-small}{hit@1}{37.4}
\setres{fresh-full}{swerank-embed-small}{hit@1-ci}{32.3--42.6}
\setres{fresh-full}{swerank-embed-small}{hit@5}{59.3}
\setres{fresh-full}{swerank-embed-small}{hit@5-ci}{52.9--65.2}
\setres{fresh-full}{swerank-embed-small}{hit@10}{65.9}
\setres{fresh-full}{swerank-embed-small}{hit@10-ci}{59.4--71.9}
\setres{fresh-full}{swerank-embed-small}{all@10}{35.0}
\setres{fresh-full}{swerank-embed-small}{all@10-ci}{29.8--40.1}
\setres{fresh-full}{swerank-embed-small}{mrr}{0.470}
\setres{fresh-full}{swerank-embed-small}{n}{580}
\setres{fresh-full}{swerank-embed-small}{cost}{--}
\setres{fresh-full}{swerank-embed-small}{sec}{0.0}
\setres{fresh-full}{website-gpu}{hit@1}{46.7}
\setres{fresh-full}{website-gpu}{hit@1-ci}{41.8--51.4}
\setres{fresh-full}{website-gpu}{hit@5}{72.1}
\setres{fresh-full}{website-gpu}{hit@5-ci}{67.3--76.6}
\setres{fresh-full}{website-gpu}{hit@10}{77.9}
\setres{fresh-full}{website-gpu}{hit@10-ci}{73.4--82.4}
\setres{fresh-full}{website-gpu}{all@10}{45.5}
\setres{fresh-full}{website-gpu}{all@10-ci}{40.0--50.9}
\setres{fresh-full}{website-gpu}{mrr}{0.575}
\setres{fresh-full}{website-gpu}{n}{580}
\setres{fresh-full}{website-gpu}{cost}{\$0.001}
\setres{fresh-full}{website-gpu}{sec}{7.5}
\setres{fresh-full}{website}{hit@1}{46.0}
\setres{fresh-full}{website}{hit@1-ci}{41.0--50.7}
\setres{fresh-full}{website}{hit@5}{68.1}
\setres{fresh-full}{website}{hit@5-ci}{62.9--72.8}
\setres{fresh-full}{website}{hit@10}{75.3}
\setres{fresh-full}{website}{hit@10-ci}{70.1--80.1}
\setres{fresh-full}{website}{all@10}{41.4}
\setres{fresh-full}{website}{all@10-ci}{35.9--46.5}
\setres{fresh-full}{website}{mrr}{0.560}
\setres{fresh-full}{website}{n}{580}
\setres{fresh-full}{website}{cost}{\$0}
\setres{fresh-full}{website}{sec}{2.0}
\setres{fresh-full}{codex-single-astra}{dhit@1}{-1.7}
\setres{fresh-full}{codex-single-astra}{wlhit@1}{89/99}
\setres{fresh-full}{codex-single-astra}{phit@1}{=1.000}
\setres{fresh-full}{codex-single-astra}{dhit@10}{-22.9}
\setres{fresh-full}{codex-single-astra}{wlhit@10}{33/166}
\setres{fresh-full}{codex-single-astra}{phit@10}{<0.001}
\setres{fresh-full}{codex-single}{dhit@1}{-12.8}
\setres{fresh-full}{codex-single}{wlhit@1}{47/121}
\setres{fresh-full}{codex-single}{phit@1}{<0.001}
\setres{fresh-full}{codex-single}{dhit@10}{-38.1}
\setres{fresh-full}{codex-single}{wlhit@10}{8/229}
\setres{fresh-full}{codex-single}{phit@10}{<0.001}
\setres{fresh-full}{gemini-single}{dhit@1}{-5.7}
\setres{fresh-full}{gemini-single}{wlhit@1}{53/77}
\setres{fresh-full}{gemini-single}{phit@1}{=0.216}
\setres{fresh-full}{gemini-single}{dhit@10}{-24.9}
\setres{fresh-full}{gemini-single}{wlhit@10}{20/126}
\setres{fresh-full}{gemini-single}{phit@10}{<0.001}
\setres{fresh-full}{open-single-gpt-oss-120b}{dhit@1}{-25.0}
\setres{fresh-full}{open-single-gpt-oss-120b}{wlhit@1}{21/166}
\setres{fresh-full}{open-single-gpt-oss-120b}{phit@1}{<0.001}
\setres{fresh-full}{open-single-gpt-oss-120b}{dhit@10}{-50.7}
\setres{fresh-full}{open-single-gpt-oss-120b}{wlhit@10}{4/298}
\setres{fresh-full}{open-single-gpt-oss-120b}{phit@10}{<0.001}
\setres{fresh-full}{open-single-qwen3-coder-next-fp8}{dhit@1}{-17.8}
\setres{fresh-full}{open-single-qwen3-coder-next-fp8}{wlhit@1}{39/142}
\setres{fresh-full}{open-single-qwen3-coder-next-fp8}{phit@1}{<0.001}
\setres{fresh-full}{open-single-qwen3-coder-next-fp8}{dhit@10}{-41.2}
\setres{fresh-full}{open-single-qwen3-coder-next-fp8}{wlhit@10}{10/249}
\setres{fresh-full}{open-single-qwen3-coder-next-fp8}{phit@10}{<0.001}
\setres{fresh-full}{swerank-embed-large}{dhit@1}{-3.5}
\setres{fresh-full}{swerank-embed-large}{wlhit@1}{71/91}
\setres{fresh-full}{swerank-embed-large}{phit@1}{=0.541}
\setres{fresh-full}{swerank-embed-large}{dhit@10}{-5.0}
\setres{fresh-full}{swerank-embed-large}{wlhit@10}{35/64}
\setres{fresh-full}{swerank-embed-large}{phit@10}{=0.046}
\setres{fresh-full}{swerank-embed-small}{dhit@1}{-9.3}
\setres{fresh-full}{swerank-embed-small}{wlhit@1}{51/105}
\setres{fresh-full}{swerank-embed-small}{phit@1}{<0.001}
\setres{fresh-full}{swerank-embed-small}{dhit@10}{-12.1}
\setres{fresh-full}{swerank-embed-small}{wlhit@10}{30/100}
\setres{fresh-full}{swerank-embed-small}{phit@10}{<0.001}
\setres{fresh-full}{website}{dhit@1}{-0.7}
\setres{fresh-full}{website}{wlhit@1}{73/77}
\setres{fresh-full}{website}{phit@1}{=1.000}
\setres{fresh-full}{website}{dhit@10}{-2.6}
\setres{fresh-full}{website}{wlhit@10}{12/27}
\setres{fresh-full}{website}{phit@10}{=0.166}
\setres{fresh-full}{codex-single-astra}{browser-dhit@1}{-1.0}
\setres{fresh-full}{codex-single-astra}{browser-wlhit@1}{93/99}
\setres{fresh-full}{codex-single-astra}{browser-phit@1}{=1.000}
\setres{fresh-full}{codex-single-astra}{browser-dhit@10}{-20.3}
\setres{fresh-full}{codex-single-astra}{browser-wlhit@10}{36/154}
\setres{fresh-full}{codex-single-astra}{browser-phit@10}{<0.001}
\setres{fresh-full}{codex-single}{browser-dhit@1}{-12.1}
\setres{fresh-full}{codex-single}{browser-wlhit@1}{61/131}
\setres{fresh-full}{codex-single}{browser-phit@1}{<0.001}
\setres{fresh-full}{codex-single}{browser-dhit@10}{-35.5}
\setres{fresh-full}{codex-single}{browser-wlhit@10}{12/218}
\setres{fresh-full}{codex-single}{browser-phit@10}{<0.001}
\setres{fresh-full}{gemini-single}{browser-dhit@1}{-5.4}
\setres{fresh-full}{gemini-single}{browser-wlhit@1}{58/81}
\setres{fresh-full}{gemini-single}{browser-phit@1}{=0.493}
\setres{fresh-full}{gemini-single}{browser-dhit@10}{-22.6}
\setres{fresh-full}{gemini-single}{browser-wlhit@10}{19/115}
\setres{fresh-full}{gemini-single}{browser-phit@10}{<0.001}
\setres{fresh-full}{open-single-gpt-oss-120b}{browser-dhit@1}{-24.3}
\setres{fresh-full}{open-single-gpt-oss-120b}{browser-wlhit@1}{33/174}
\setres{fresh-full}{open-single-gpt-oss-120b}{browser-phit@1}{<0.001}
\setres{fresh-full}{open-single-gpt-oss-120b}{browser-dhit@10}{-48.1}
\setres{fresh-full}{open-single-gpt-oss-120b}{browser-wlhit@10}{4/283}
\setres{fresh-full}{open-single-gpt-oss-120b}{browser-phit@10}{<0.001}
\setres{fresh-full}{open-single-qwen3-coder-next-fp8}{browser-dhit@1}{-17.1}
\setres{fresh-full}{open-single-qwen3-coder-next-fp8}{browser-wlhit@1}{49/148}
\setres{fresh-full}{open-single-qwen3-coder-next-fp8}{browser-phit@1}{<0.001}
\setres{fresh-full}{open-single-qwen3-coder-next-fp8}{browser-dhit@10}{-38.6}
\setres{fresh-full}{open-single-qwen3-coder-next-fp8}{browser-wlhit@10}{13/237}
\setres{fresh-full}{open-single-qwen3-coder-next-fp8}{browser-phit@10}{<0.001}
\setres{fresh-full}{swerank-embed-large}{browser-dhit@1}{-2.8}
\setres{fresh-full}{swerank-embed-large}{browser-wlhit@1}{75/91}
\setres{fresh-full}{swerank-embed-large}{browser-phit@1}{=1.000}
\setres{fresh-full}{swerank-embed-large}{browser-dhit@10}{-2.4}
\setres{fresh-full}{swerank-embed-large}{browser-wlhit@10}{46/60}
\setres{fresh-full}{swerank-embed-large}{browser-phit@10}{=1.000}
\setres{fresh-full}{swerank-embed-small}{browser-dhit@1}{-8.6}
\setres{fresh-full}{swerank-embed-small}{browser-wlhit@1}{55/105}
\setres{fresh-full}{swerank-embed-small}{browser-phit@1}{=0.002}
\setres{fresh-full}{swerank-embed-small}{browser-dhit@10}{-9.5}
\setres{fresh-full}{swerank-embed-small}{browser-wlhit@10}{38/93}
\setres{fresh-full}{swerank-embed-small}{browser-phit@10}{<0.001}
\setres{fresh-full}{website-gpu}{browser-dhit@1}{+0.7}
\setres{fresh-full}{website-gpu}{browser-wlhit@1}{77/73}
\setres{fresh-full}{website-gpu}{browser-phit@1}{=1.000}
\setres{fresh-full}{website-gpu}{browser-dhit@10}{+2.6}
\setres{fresh-full}{website-gpu}{browser-wlhit@10}{27/12}
\setres{fresh-full}{website-gpu}{browser-phit@10}{=0.237}
\setres{fresh-full}{codex-single-astra}{s-function}{82.1}
\setres{fresh-full}{codex-single-astra}{sn-function}{246}
\setres{fresh-full}{codex-single-astra}{s-file}{51.5}
\setres{fresh-full}{codex-single-astra}{sn-file}{68}
\setres{fresh-full}{codex-single-astra}{s-trace}{0.0}
\setres{fresh-full}{codex-single-astra}{sn-trace}{6}
\setres{fresh-full}{codex-single-astra}{s-none}{31.5}
\setres{fresh-full}{codex-single-astra}{sn-none}{260}
\setres{fresh-full}{codex-single}{s-function}{73.6}
\setres{fresh-full}{codex-single}{sn-function}{246}
\setres{fresh-full}{codex-single}{s-file}{32.4}
\setres{fresh-full}{codex-single}{sn-file}{68}
\setres{fresh-full}{codex-single}{s-trace}{0.0}
\setres{fresh-full}{codex-single}{sn-trace}{6}
\setres{fresh-full}{codex-single}{s-none}{10.8}
\setres{fresh-full}{codex-single}{sn-none}{260}
\setres{fresh-full}{gemini-single}{s-function}{83.9}
\setres{fresh-full}{gemini-single}{sn-function}{186}
\setres{fresh-full}{gemini-single}{s-file}{43.1}
\setres{fresh-full}{gemini-single}{sn-file}{58}
\setres{fresh-full}{gemini-single}{s-trace}{0.0}
\setres{fresh-full}{gemini-single}{sn-trace}{6}
\setres{fresh-full}{gemini-single}{s-none}{24.6}
\setres{fresh-full}{gemini-single}{sn-none}{175}
\setres{fresh-full}{open-single-gpt-oss-120b}{s-function}{53.2}
\setres{fresh-full}{open-single-gpt-oss-120b}{sn-function}{246}
\setres{fresh-full}{open-single-gpt-oss-120b}{s-file}{14.7}
\setres{fresh-full}{open-single-gpt-oss-120b}{sn-file}{68}
\setres{fresh-full}{open-single-gpt-oss-120b}{s-trace}{0.0}
\setres{fresh-full}{open-single-gpt-oss-120b}{sn-trace}{6}
\setres{fresh-full}{open-single-gpt-oss-120b}{s-none}{6.5}
\setres{fresh-full}{open-single-gpt-oss-120b}{sn-none}{260}
\setres{fresh-full}{open-single-qwen3-coder-next-fp8}{s-function}{69.5}
\setres{fresh-full}{open-single-qwen3-coder-next-fp8}{sn-function}{246}
\setres{fresh-full}{open-single-qwen3-coder-next-fp8}{s-file}{19.1}
\setres{fresh-full}{open-single-qwen3-coder-next-fp8}{sn-file}{68}
\setres{fresh-full}{open-single-qwen3-coder-next-fp8}{s-trace}{0.0}
\setres{fresh-full}{open-single-qwen3-coder-next-fp8}{sn-trace}{6}
\setres{fresh-full}{open-single-qwen3-coder-next-fp8}{s-none}{11.2}
\setres{fresh-full}{open-single-qwen3-coder-next-fp8}{sn-none}{260}
\setres{fresh-full}{swerank-embed-large}{s-function}{84.5}
\setres{fresh-full}{swerank-embed-large}{sn-function}{246}
\setres{fresh-full}{swerank-embed-large}{s-file}{73.5}
\setres{fresh-full}{swerank-embed-large}{sn-file}{68}
\setres{fresh-full}{swerank-embed-large}{s-trace}{16.7}
\setres{fresh-full}{swerank-embed-large}{sn-trace}{6}
\setres{fresh-full}{swerank-embed-large}{s-none}{63.1}
\setres{fresh-full}{swerank-embed-large}{sn-none}{260}
\setres{fresh-full}{swerank-embed-small}{s-function}{77.2}
\setres{fresh-full}{swerank-embed-small}{sn-function}{246}
\setres{fresh-full}{swerank-embed-small}{s-file}{67.7}
\setres{fresh-full}{swerank-embed-small}{sn-file}{68}
\setres{fresh-full}{swerank-embed-small}{s-trace}{16.7}
\setres{fresh-full}{swerank-embed-small}{sn-trace}{6}
\setres{fresh-full}{swerank-embed-small}{s-none}{55.8}
\setres{fresh-full}{swerank-embed-small}{sn-none}{260}
\setres{fresh-full}{website-gpu}{s-function}{95.1}
\setres{fresh-full}{website-gpu}{sn-function}{246}
\setres{fresh-full}{website-gpu}{s-file}{83.8}
\setres{fresh-full}{website-gpu}{sn-file}{68}
\setres{fresh-full}{website-gpu}{s-trace}{16.7}
\setres{fresh-full}{website-gpu}{sn-trace}{6}
\setres{fresh-full}{website-gpu}{s-none}{61.5}
\setres{fresh-full}{website-gpu}{sn-none}{260}
\setres{fresh-full}{website}{s-function}{93.9}
\setres{fresh-full}{website}{sn-function}{246}
\setres{fresh-full}{website}{s-file}{86.8}
\setres{fresh-full}{website}{sn-file}{68}
\setres{fresh-full}{website}{s-trace}{0.0}
\setres{fresh-full}{website}{sn-trace}{6}
\setres{fresh-full}{website}{s-none}{56.5}
\setres{fresh-full}{website}{sn-none}{260}
\setres{fresh-full}{codex-single-astra}{esc-escalated}{58.6}
\setres{fresh-full}{codex-single-astra}{esc-hit@10}{64.5}
\setres{fresh-full}{codex-single-astra}{esc-random_hit@10}{64.5}
\setres{fresh-full}{codex-single}{esc-escalated}{58.6}
\setres{fresh-full}{codex-single}{esc-hit@10}{52.2}
\setres{fresh-full}{codex-single}{esc-random_hit@10}{55.6}
\setres{fresh-full}{gemini-single}{esc-escalated}{56.2}
\setres{fresh-full}{gemini-single}{esc-hit@10}{61.4}
\setres{fresh-full}{gemini-single}{esc-random_hit@10}{63.6}
\setres{fresh-full}{open-single-gpt-oss-120b}{esc-escalated}{58.6}
\setres{fresh-full}{open-single-gpt-oss-120b}{esc-hit@10}{46.7}
\setres{fresh-full}{open-single-gpt-oss-120b}{esc-random_hit@10}{48.2}
\setres{fresh-full}{open-single-qwen3-coder-next-fp8}{esc-escalated}{58.6}
\setres{fresh-full}{open-single-qwen3-coder-next-fp8}{esc-hit@10}{51.4}
\setres{fresh-full}{open-single-qwen3-coder-next-fp8}{esc-random_hit@10}{53.8}
\setres{fresh-full}{swerank-embed-large}{esc-escalated}{58.6}
\setres{fresh-full}{swerank-embed-large}{esc-hit@10}{78.3}
\setres{fresh-full}{swerank-embed-large}{esc-random_hit@10}{75.0}
\setres{fresh-full}{swerank-embed-small}{esc-escalated}{58.6}
\setres{fresh-full}{swerank-embed-small}{esc-hit@10}{72.8}
\setres{fresh-full}{swerank-embed-small}{esc-random_hit@10}{70.9}
\setres{fresh-full}{website}{esc-escalated}{58.6}
\setres{fresh-full}{website}{esc-hit@10}{75.3}
\setres{fresh-full}{website}{esc-random_hit@10}{76.4}
\setres{fresh-agents}{claude-agent}{hit@1}{77.8}
\setres{fresh-agents}{claude-agent}{hit@1-ci}{71.4--83.6}
\setres{fresh-agents}{claude-agent}{hit@5}{89.9}
\setres{fresh-agents}{claude-agent}{hit@5-ci}{84.9--94.2}
\setres{fresh-agents}{claude-agent}{hit@10}{90.9}
\setres{fresh-agents}{claude-agent}{hit@10-ci}{86.2--95.1}
\setres{fresh-agents}{claude-agent}{all@10}{49.5}
\setres{fresh-agents}{claude-agent}{all@10-ci}{42.4--56.7}
\setres{fresh-agents}{claude-agent}{mrr}{0.828}
\setres{fresh-agents}{claude-agent}{n}{198}
\setres{fresh-agents}{claude-agent}{cost}{\$0.20}
\setres{fresh-agents}{claude-agent}{sec}{27.2}
\setres{fresh-agents}{claude-single-sonnet}{hit@1}{50.0}
\setres{fresh-agents}{claude-single-sonnet}{hit@1-ci}{50.0--50.0}
\setres{fresh-agents}{claude-single-sonnet}{hit@5}{100.0}
\setres{fresh-agents}{claude-single-sonnet}{hit@5-ci}{100.0--100.0}
\setres{fresh-agents}{claude-single-sonnet}{hit@10}{100.0}
\setres{fresh-agents}{claude-single-sonnet}{hit@10-ci}{100.0--100.0}
\setres{fresh-agents}{claude-single-sonnet}{all@10}{50.0}
\setres{fresh-agents}{claude-single-sonnet}{all@10-ci}{50.0--50.0}
\setres{fresh-agents}{claude-single-sonnet}{mrr}{0.667}
\setres{fresh-agents}{claude-single-sonnet}{n}{2}
\setres{fresh-agents}{claude-single-sonnet}{cost}{\$0.10}
\setres{fresh-agents}{claude-single-sonnet}{sec}{8.6}
\setres{fresh-agents}{codex-agent-astra}{hit@1}{80.3}
\setres{fresh-agents}{codex-agent-astra}{hit@1-ci}{74.6--85.5}
\setres{fresh-agents}{codex-agent-astra}{hit@5}{86.4}
\setres{fresh-agents}{codex-agent-astra}{hit@5-ci}{80.9--91.2}
\setres{fresh-agents}{codex-agent-astra}{hit@10}{86.4}
\setres{fresh-agents}{codex-agent-astra}{hit@10-ci}{80.9--91.2}
\setres{fresh-agents}{codex-agent-astra}{all@10}{41.4}
\setres{fresh-agents}{codex-agent-astra}{all@10-ci}{34.5--48.5}
\setres{fresh-agents}{codex-agent-astra}{mrr}{0.830}
\setres{fresh-agents}{codex-agent-astra}{n}{198}
\setres{fresh-agents}{codex-agent-astra}{cost}{\$0.29}
\setres{fresh-agents}{codex-agent-astra}{sec}{24.0}
\setres{fresh-agents}{codex-agent}{hit@1}{67.2}
\setres{fresh-agents}{codex-agent}{hit@1-ci}{60.2--73.6}
\setres{fresh-agents}{codex-agent}{hit@5}{75.8}
\setres{fresh-agents}{codex-agent}{hit@5-ci}{70.0--81.2}
\setres{fresh-agents}{codex-agent}{hit@10}{75.8}
\setres{fresh-agents}{codex-agent}{hit@10-ci}{70.0--81.2}
\setres{fresh-agents}{codex-agent}{all@10}{34.3}
\setres{fresh-agents}{codex-agent}{all@10-ci}{27.8--41.1}
\setres{fresh-agents}{codex-agent}{mrr}{0.710}
\setres{fresh-agents}{codex-agent}{n}{198}
\setres{fresh-agents}{codex-agent}{cost}{\$0.004}
\setres{fresh-agents}{codex-agent}{sec}{20.1}
\setres{fresh-agents}{codex-single-astra}{hit@1}{43.4}
\setres{fresh-agents}{codex-single-astra}{hit@1-ci}{34.8--52.2}
\setres{fresh-agents}{codex-single-astra}{hit@5}{56.1}
\setres{fresh-agents}{codex-single-astra}{hit@5-ci}{47.8--64.4}
\setres{fresh-agents}{codex-single-astra}{hit@10}{56.1}
\setres{fresh-agents}{codex-single-astra}{hit@10-ci}{47.8--64.4}
\setres{fresh-agents}{codex-single-astra}{all@10}{21.7}
\setres{fresh-agents}{codex-single-astra}{all@10-ci}{15.5--28.6}
\setres{fresh-agents}{codex-single-astra}{mrr}{0.489}
\setres{fresh-agents}{codex-single-astra}{n}{198}
\setres{fresh-agents}{codex-single-astra}{cost}{\$0.51}
\setres{fresh-agents}{codex-single-astra}{sec}{14.8}
\setres{fresh-agents}{codex-single}{hit@1}{30.3}
\setres{fresh-agents}{codex-single}{hit@1-ci}{22.2--38.7}
\setres{fresh-agents}{codex-single}{hit@5}{37.9}
\setres{fresh-agents}{codex-single}{hit@5-ci}{29.6--46.0}
\setres{fresh-agents}{codex-single}{hit@10}{37.9}
\setres{fresh-agents}{codex-single}{hit@10-ci}{29.6--46.0}
\setres{fresh-agents}{codex-single}{all@10}{15.7}
\setres{fresh-agents}{codex-single}{all@10-ci}{10.5--21.3}
\setres{fresh-agents}{codex-single}{mrr}{0.338}
\setres{fresh-agents}{codex-single}{n}{198}
\setres{fresh-agents}{codex-single}{cost}{\$0.005}
\setres{fresh-agents}{codex-single}{sec}{9.7}
\setres{fresh-agents}{gemini-single}{hit@1}{36.4}
\setres{fresh-agents}{gemini-single}{hit@1-ci}{28.0--44.9}
\setres{fresh-agents}{gemini-single}{hit@5}{49.5}
\setres{fresh-agents}{gemini-single}{hit@5-ci}{41.2--57.9}
\setres{fresh-agents}{gemini-single}{hit@10}{51.5}
\setres{fresh-agents}{gemini-single}{hit@10-ci}{43.0--60.3}
\setres{fresh-agents}{gemini-single}{all@10}{22.2}
\setres{fresh-agents}{gemini-single}{all@10-ci}{16.2--29.1}
\setres{fresh-agents}{gemini-single}{mrr}{0.417}
\setres{fresh-agents}{gemini-single}{n}{198}
\setres{fresh-agents}{gemini-single}{cost}{\$0.11}
\setres{fresh-agents}{gemini-single}{sec}{38.0}
\setres{fresh-agents}{open-agent-qwen3-coder-next-fp8}{hit@1}{52.5}
\setres{fresh-agents}{open-agent-qwen3-coder-next-fp8}{hit@1-ci}{44.0--61.3}
\setres{fresh-agents}{open-agent-qwen3-coder-next-fp8}{hit@5}{60.1}
\setres{fresh-agents}{open-agent-qwen3-coder-next-fp8}{hit@5-ci}{52.2--68.1}
\setres{fresh-agents}{open-agent-qwen3-coder-next-fp8}{hit@10}{60.6}
\setres{fresh-agents}{open-agent-qwen3-coder-next-fp8}{hit@10-ci}{52.7--68.7}
\setres{fresh-agents}{open-agent-qwen3-coder-next-fp8}{all@10}{29.8}
\setres{fresh-agents}{open-agent-qwen3-coder-next-fp8}{all@10-ci}{22.8--37.2}
\setres{fresh-agents}{open-agent-qwen3-coder-next-fp8}{mrr}{0.557}
\setres{fresh-agents}{open-agent-qwen3-coder-next-fp8}{n}{198}
\setres{fresh-agents}{open-agent-qwen3-coder-next-fp8}{cost}{\$0.028}
\setres{fresh-agents}{open-agent-qwen3-coder-next-fp8}{sec}{20.4}
\setres{fresh-agents}{open-single-gpt-oss-120b}{hit@1}{20.7}
\setres{fresh-agents}{open-single-gpt-oss-120b}{hit@1-ci}{13.8--28.0}
\setres{fresh-agents}{open-single-gpt-oss-120b}{hit@5}{24.7}
\setres{fresh-agents}{open-single-gpt-oss-120b}{hit@5-ci}{17.0--32.6}
\setres{fresh-agents}{open-single-gpt-oss-120b}{hit@10}{24.7}
\setres{fresh-agents}{open-single-gpt-oss-120b}{hit@10-ci}{17.0--32.6}
\setres{fresh-agents}{open-single-gpt-oss-120b}{all@10}{10.6}
\setres{fresh-agents}{open-single-gpt-oss-120b}{all@10-ci}{6.1--15.6}
\setres{fresh-agents}{open-single-gpt-oss-120b}{mrr}{0.224}
\setres{fresh-agents}{open-single-gpt-oss-120b}{n}{198}
\setres{fresh-agents}{open-single-gpt-oss-120b}{cost}{\$0.006}
\setres{fresh-agents}{open-single-gpt-oss-120b}{sec}{4.4}
\setres{fresh-agents}{open-single-qwen3-coder-next-fp8}{hit@1}{24.2}
\setres{fresh-agents}{open-single-qwen3-coder-next-fp8}{hit@1-ci}{16.8--32.3}
\setres{fresh-agents}{open-single-qwen3-coder-next-fp8}{hit@5}{32.3}
\setres{fresh-agents}{open-single-qwen3-coder-next-fp8}{hit@5-ci}{24.7--40.4}
\setres{fresh-agents}{open-single-qwen3-coder-next-fp8}{hit@10}{33.3}
\setres{fresh-agents}{open-single-qwen3-coder-next-fp8}{hit@10-ci}{25.8--41.3}
\setres{fresh-agents}{open-single-qwen3-coder-next-fp8}{all@10}{11.1}
\setres{fresh-agents}{open-single-qwen3-coder-next-fp8}{all@10-ci}{6.5--16.3}
\setres{fresh-agents}{open-single-qwen3-coder-next-fp8}{mrr}{0.277}
\setres{fresh-agents}{open-single-qwen3-coder-next-fp8}{n}{198}
\setres{fresh-agents}{open-single-qwen3-coder-next-fp8}{cost}{\$0.007}
\setres{fresh-agents}{open-single-qwen3-coder-next-fp8}{sec}{3.0}
\setres{fresh-agents}{swerank-embed-large}{hit@1}{42.9}
\setres{fresh-agents}{swerank-embed-large}{hit@1-ci}{35.4--50.5}
\setres{fresh-agents}{swerank-embed-large}{hit@5}{64.6}
\setres{fresh-agents}{swerank-embed-large}{hit@5-ci}{57.6--71.7}
\setres{fresh-agents}{swerank-embed-large}{hit@10}{70.2}
\setres{fresh-agents}{swerank-embed-large}{hit@10-ci}{63.4--76.9}
\setres{fresh-agents}{swerank-embed-large}{all@10}{37.4}
\setres{fresh-agents}{swerank-embed-large}{all@10-ci}{30.2--44.9}
\setres{fresh-agents}{swerank-embed-large}{mrr}{0.520}
\setres{fresh-agents}{swerank-embed-large}{n}{198}
\setres{fresh-agents}{swerank-embed-large}{cost}{--}
\setres{fresh-agents}{swerank-embed-large}{sec}{0.0}
\setres{fresh-agents}{swerank-embed-small}{hit@1}{39.4}
\setres{fresh-agents}{swerank-embed-small}{hit@1-ci}{31.6--47.7}
\setres{fresh-agents}{swerank-embed-small}{hit@5}{55.1}
\setres{fresh-agents}{swerank-embed-small}{hit@5-ci}{46.8--63.3}
\setres{fresh-agents}{swerank-embed-small}{hit@10}{62.6}
\setres{fresh-agents}{swerank-embed-small}{hit@10-ci}{54.6--70.2}
\setres{fresh-agents}{swerank-embed-small}{all@10}{32.8}
\setres{fresh-agents}{swerank-embed-small}{all@10-ci}{26.7--39.1}
\setres{fresh-agents}{swerank-embed-small}{mrr}{0.456}
\setres{fresh-agents}{swerank-embed-small}{n}{198}
\setres{fresh-agents}{swerank-embed-small}{cost}{--}
\setres{fresh-agents}{swerank-embed-small}{sec}{0.0}
\setres{fresh-agents}{website-gpu}{hit@1}{45.5}
\setres{fresh-agents}{website-gpu}{hit@1-ci}{38.3--52.8}
\setres{fresh-agents}{website-gpu}{hit@5}{66.7}
\setres{fresh-agents}{website-gpu}{hit@5-ci}{60.1--73.0}
\setres{fresh-agents}{website-gpu}{hit@10}{74.2}
\setres{fresh-agents}{website-gpu}{hit@10-ci}{68.3--79.9}
\setres{fresh-agents}{website-gpu}{all@10}{43.4}
\setres{fresh-agents}{website-gpu}{all@10-ci}{35.9--50.6}
\setres{fresh-agents}{website-gpu}{mrr}{0.540}
\setres{fresh-agents}{website-gpu}{n}{198}
\setres{fresh-agents}{website-gpu}{cost}{\$0.001}
\setres{fresh-agents}{website-gpu}{sec}{10.8}
\setres{fresh-agents}{website}{hit@1}{41.9}
\setres{fresh-agents}{website}{hit@1-ci}{34.3--49.5}
\setres{fresh-agents}{website}{hit@5}{65.2}
\setres{fresh-agents}{website}{hit@5-ci}{58.4--72.0}
\setres{fresh-agents}{website}{hit@10}{73.7}
\setres{fresh-agents}{website}{hit@10-ci}{68.4--79.4}
\setres{fresh-agents}{website}{all@10}{39.4}
\setres{fresh-agents}{website}{all@10-ci}{32.6--46.5}
\setres{fresh-agents}{website}{mrr}{0.520}
\setres{fresh-agents}{website}{n}{198}
\setres{fresh-agents}{website}{cost}{\$0}
\setres{fresh-agents}{website}{sec}{4.8}
\setres{fresh-agents}{claude-agent}{dhit@1}{+32.3}
\setres{fresh-agents}{claude-agent}{wlhit@1}{78/14}
\setres{fresh-agents}{claude-agent}{phit@1}{<0.001}
\setres{fresh-agents}{claude-agent}{dhit@10}{+16.7}
\setres{fresh-agents}{claude-agent}{wlhit@10}{41/8}
\setres{fresh-agents}{claude-agent}{phit@10}{<0.001}
\setres{fresh-agents}{claude-single-sonnet}{dhit@1}{+0.0}
\setres{fresh-agents}{claude-single-sonnet}{wlhit@1}{0/0}
\setres{fresh-agents}{claude-single-sonnet}{phit@1}{=1.000}
\setres{fresh-agents}{claude-single-sonnet}{dhit@10}{+50.0}
\setres{fresh-agents}{claude-single-sonnet}{wlhit@10}{1/0}
\setres{fresh-agents}{claude-single-sonnet}{phit@10}{=1.000}
\setres{fresh-agents}{codex-agent-astra}{dhit@1}{+34.8}
\setres{fresh-agents}{codex-agent-astra}{wlhit@1}{80/11}
\setres{fresh-agents}{codex-agent-astra}{phit@1}{<0.001}
\setres{fresh-agents}{codex-agent-astra}{dhit@10}{+12.1}
\setres{fresh-agents}{codex-agent-astra}{wlhit@10}{38/14}
\setres{fresh-agents}{codex-agent-astra}{phit@10}{=0.038}
\setres{fresh-agents}{codex-agent}{dhit@1}{+21.7}
\setres{fresh-agents}{codex-agent}{wlhit@1}{62/19}
\setres{fresh-agents}{codex-agent}{phit@1}{<0.001}
\setres{fresh-agents}{codex-agent}{dhit@10}{+1.5}
\setres{fresh-agents}{codex-agent}{wlhit@10}{26/23}
\setres{fresh-agents}{codex-agent}{phit@10}{=1.000}
\setres{fresh-agents}{codex-single-astra}{dhit@1}{-2.0}
\setres{fresh-agents}{codex-single-astra}{wlhit@1}{35/39}
\setres{fresh-agents}{codex-single-astra}{phit@1}{=1.000}
\setres{fresh-agents}{codex-single-astra}{dhit@10}{-18.2}
\setres{fresh-agents}{codex-single-astra}{wlhit@10}{13/49}
\setres{fresh-agents}{codex-single-astra}{phit@10}{<0.001}
\setres{fresh-agents}{codex-single}{dhit@1}{-15.2}
\setres{fresh-agents}{codex-single}{wlhit@1}{15/45}
\setres{fresh-agents}{codex-single}{phit@1}{=0.004}
\setres{fresh-agents}{codex-single}{dhit@10}{-36.4}
\setres{fresh-agents}{codex-single}{wlhit@10}{4/76}
\setres{fresh-agents}{codex-single}{phit@10}{<0.001}
\setres{fresh-agents}{gemini-single}{dhit@1}{-9.1}
\setres{fresh-agents}{gemini-single}{wlhit@1}{26/44}
\setres{fresh-agents}{gemini-single}{phit@1}{=1.000}
\setres{fresh-agents}{gemini-single}{dhit@10}{-22.7}
\setres{fresh-agents}{gemini-single}{wlhit@10}{13/58}
\setres{fresh-agents}{gemini-single}{phit@10}{<0.001}
\setres{fresh-agents}{open-agent-qwen3-coder-next-fp8}{dhit@1}{+7.1}
\setres{fresh-agents}{open-agent-qwen3-coder-next-fp8}{wlhit@1}{40/26}
\setres{fresh-agents}{open-agent-qwen3-coder-next-fp8}{phit@1}{=1.000}
\setres{fresh-agents}{open-agent-qwen3-coder-next-fp8}{dhit@10}{-13.6}
\setres{fresh-agents}{open-agent-qwen3-coder-next-fp8}{wlhit@10}{19/46}
\setres{fresh-agents}{open-agent-qwen3-coder-next-fp8}{phit@10}{=0.036}
\setres{fresh-agents}{open-single-gpt-oss-120b}{dhit@1}{-24.8}
\setres{fresh-agents}{open-single-gpt-oss-120b}{wlhit@1}{7/56}
\setres{fresh-agents}{open-single-gpt-oss-120b}{phit@1}{<0.001}
\setres{fresh-agents}{open-single-gpt-oss-120b}{dhit@10}{-49.5}
\setres{fresh-agents}{open-single-gpt-oss-120b}{wlhit@10}{1/99}
\setres{fresh-agents}{open-single-gpt-oss-120b}{phit@10}{<0.001}
\setres{fresh-agents}{open-single-qwen3-coder-next-fp8}{dhit@1}{-21.2}
\setres{fresh-agents}{open-single-qwen3-coder-next-fp8}{wlhit@1}{14/56}
\setres{fresh-agents}{open-single-qwen3-coder-next-fp8}{phit@1}{<0.001}
\setres{fresh-agents}{open-single-qwen3-coder-next-fp8}{dhit@10}{-40.9}
\setres{fresh-agents}{open-single-qwen3-coder-next-fp8}{wlhit@10}{5/86}
\setres{fresh-agents}{open-single-qwen3-coder-next-fp8}{phit@10}{<0.001}
\setres{fresh-agents}{swerank-embed-large}{dhit@1}{-2.5}
\setres{fresh-agents}{swerank-embed-large}{wlhit@1}{28/33}
\setres{fresh-agents}{swerank-embed-large}{phit@1}{=1.000}
\setres{fresh-agents}{swerank-embed-large}{dhit@10}{-4.0}
\setres{fresh-agents}{swerank-embed-large}{wlhit@10}{18/26}
\setres{fresh-agents}{swerank-embed-large}{phit@10}{=1.000}
\setres{fresh-agents}{swerank-embed-small}{dhit@1}{-6.1}
\setres{fresh-agents}{swerank-embed-small}{wlhit@1}{21/33}
\setres{fresh-agents}{swerank-embed-small}{phit@1}{=1.000}
\setres{fresh-agents}{swerank-embed-small}{dhit@10}{-11.6}
\setres{fresh-agents}{swerank-embed-small}{wlhit@10}{16/39}
\setres{fresh-agents}{swerank-embed-small}{phit@10}{=0.084}
\setres{fresh-agents}{website}{dhit@1}{-3.5}
\setres{fresh-agents}{website}{wlhit@1}{21/28}
\setres{fresh-agents}{website}{phit@1}{=1.000}
\setres{fresh-agents}{website}{dhit@10}{-0.5}
\setres{fresh-agents}{website}{wlhit@10}{7/8}
\setres{fresh-agents}{website}{phit@10}{=1.000}
\setres{fresh-agents}{claude-agent}{browser-dhit@1}{+35.9}
\setres{fresh-agents}{claude-agent}{browser-wlhit@1}{80/9}
\setres{fresh-agents}{claude-agent}{browser-phit@1}{<0.001}
\setres{fresh-agents}{claude-agent}{browser-dhit@10}{+17.2}
\setres{fresh-agents}{claude-agent}{browser-wlhit@10}{43/9}
\setres{fresh-agents}{claude-agent}{browser-phit@10}{<0.001}
\setres{fresh-agents}{claude-single-sonnet}{browser-dhit@1}{+50.0}
\setres{fresh-agents}{claude-single-sonnet}{browser-wlhit@1}{1/0}
\setres{fresh-agents}{claude-single-sonnet}{browser-phit@1}{=1.000}
\setres{fresh-agents}{claude-single-sonnet}{browser-dhit@10}{+50.0}
\setres{fresh-agents}{claude-single-sonnet}{browser-wlhit@10}{1/0}
\setres{fresh-agents}{claude-single-sonnet}{browser-phit@10}{=1.000}
\setres{fresh-agents}{codex-agent-astra}{browser-dhit@1}{+38.4}
\setres{fresh-agents}{codex-agent-astra}{browser-wlhit@1}{86/10}
\setres{fresh-agents}{codex-agent-astra}{browser-phit@1}{<0.001}
\setres{fresh-agents}{codex-agent-astra}{browser-dhit@10}{+12.6}
\setres{fresh-agents}{codex-agent-astra}{browser-wlhit@10}{39/14}
\setres{fresh-agents}{codex-agent-astra}{browser-phit@10}{=0.030}
\setres{fresh-agents}{codex-agent}{browser-dhit@1}{+25.2}
\setres{fresh-agents}{codex-agent}{browser-wlhit@1}{65/15}
\setres{fresh-agents}{codex-agent}{browser-phit@1}{<0.001}
\setres{fresh-agents}{codex-agent}{browser-dhit@10}{+2.0}
\setres{fresh-agents}{codex-agent}{browser-wlhit@10}{27/23}
\setres{fresh-agents}{codex-agent}{browser-phit@10}{=1.000}
\setres{fresh-agents}{codex-single-astra}{browser-dhit@1}{+1.5}
\setres{fresh-agents}{codex-single-astra}{browser-wlhit@1}{34/31}
\setres{fresh-agents}{codex-single-astra}{browser-phit@1}{=1.000}
\setres{fresh-agents}{codex-single-astra}{browser-dhit@10}{-17.7}
\setres{fresh-agents}{codex-single-astra}{browser-wlhit@10}{14/49}
\setres{fresh-agents}{codex-single-astra}{browser-phit@10}{<0.001}
\setres{fresh-agents}{codex-single}{browser-dhit@1}{-11.6}
\setres{fresh-agents}{codex-single}{browser-wlhit@1}{17/40}
\setres{fresh-agents}{codex-single}{browser-phit@1}{=0.109}
\setres{fresh-agents}{codex-single}{browser-dhit@10}{-35.9}
\setres{fresh-agents}{codex-single}{browser-wlhit@10}{4/75}
\setres{fresh-agents}{codex-single}{browser-phit@10}{<0.001}
\setres{fresh-agents}{gemini-single}{browser-dhit@1}{-5.6}
\setres{fresh-agents}{gemini-single}{browser-wlhit@1}{23/34}
\setres{fresh-agents}{gemini-single}{browser-phit@1}{=1.000}
\setres{fresh-agents}{gemini-single}{browser-dhit@10}{-22.2}
\setres{fresh-agents}{gemini-single}{browser-wlhit@10}{13/57}
\setres{fresh-agents}{gemini-single}{browser-phit@10}{<0.001}
\setres{fresh-agents}{open-agent-qwen3-coder-next-fp8}{browser-dhit@1}{+10.6}
\setres{fresh-agents}{open-agent-qwen3-coder-next-fp8}{browser-wlhit@1}{48/27}
\setres{fresh-agents}{open-agent-qwen3-coder-next-fp8}{browser-phit@1}{=0.568}
\setres{fresh-agents}{open-agent-qwen3-coder-next-fp8}{browser-dhit@10}{-13.1}
\setres{fresh-agents}{open-agent-qwen3-coder-next-fp8}{browser-wlhit@10}{22/48}
\setres{fresh-agents}{open-agent-qwen3-coder-next-fp8}{browser-phit@10}{=0.087}
\setres{fresh-agents}{open-single-gpt-oss-120b}{browser-dhit@1}{-21.2}
\setres{fresh-agents}{open-single-gpt-oss-120b}{browser-wlhit@1}{7/49}
\setres{fresh-agents}{open-single-gpt-oss-120b}{browser-phit@1}{<0.001}
\setres{fresh-agents}{open-single-gpt-oss-120b}{browser-dhit@10}{-49.0}
\setres{fresh-agents}{open-single-gpt-oss-120b}{browser-wlhit@10}{1/98}
\setres{fresh-agents}{open-single-gpt-oss-120b}{browser-phit@10}{<0.001}
\setres{fresh-agents}{open-single-qwen3-coder-next-fp8}{browser-dhit@1}{-17.7}
\setres{fresh-agents}{open-single-qwen3-coder-next-fp8}{browser-wlhit@1}{15/50}
\setres{fresh-agents}{open-single-qwen3-coder-next-fp8}{browser-phit@1}{<0.001}
\setres{fresh-agents}{open-single-qwen3-coder-next-fp8}{browser-dhit@10}{-40.4}
\setres{fresh-agents}{open-single-qwen3-coder-next-fp8}{browser-wlhit@10}{4/84}
\setres{fresh-agents}{open-single-qwen3-coder-next-fp8}{browser-phit@10}{<0.001}
\setres{fresh-agents}{swerank-embed-large}{browser-dhit@1}{+1.0}
\setres{fresh-agents}{swerank-embed-large}{browser-wlhit@1}{30/28}
\setres{fresh-agents}{swerank-embed-large}{browser-phit@1}{=1.000}
\setres{fresh-agents}{swerank-embed-large}{browser-dhit@10}{-3.5}
\setres{fresh-agents}{swerank-embed-large}{browser-wlhit@10}{20/27}
\setres{fresh-agents}{swerank-embed-large}{browser-phit@10}{=1.000}
\setres{fresh-agents}{swerank-embed-small}{browser-dhit@1}{-2.5}
\setres{fresh-agents}{swerank-embed-small}{browser-wlhit@1}{27/32}
\setres{fresh-agents}{swerank-embed-small}{browser-phit@1}{=1.000}
\setres{fresh-agents}{swerank-embed-small}{browser-dhit@10}{-11.1}
\setres{fresh-agents}{swerank-embed-small}{browser-wlhit@10}{16/38}
\setres{fresh-agents}{swerank-embed-small}{browser-phit@10}{=0.125}
\setres{fresh-agents}{website-gpu}{browser-dhit@1}{+3.5}
\setres{fresh-agents}{website-gpu}{browser-wlhit@1}{28/21}
\setres{fresh-agents}{website-gpu}{browser-phit@1}{=1.000}
\setres{fresh-agents}{website-gpu}{browser-dhit@10}{+0.5}
\setres{fresh-agents}{website-gpu}{browser-wlhit@10}{8/7}
\setres{fresh-agents}{website-gpu}{browser-phit@10}{=1.000}
\setres{fresh-agents}{claude-agent}{s-function}{95.3}
\setres{fresh-agents}{claude-agent}{sn-function}{64}
\setres{fresh-agents}{claude-agent}{s-file}{100.0}
\setres{fresh-agents}{claude-agent}{sn-file}{14}
\setres{fresh-agents}{claude-agent}{s-trace}{66.7}
\setres{fresh-agents}{claude-agent}{sn-trace}{3}
\setres{fresh-agents}{claude-agent}{s-none}{88.0}
\setres{fresh-agents}{claude-agent}{sn-none}{117}
\setres{fresh-agents}{claude-single-sonnet}{s-trace}{100.0}
\setres{fresh-agents}{claude-single-sonnet}{sn-trace}{1}
\setres{fresh-agents}{claude-single-sonnet}{s-none}{100.0}
\setres{fresh-agents}{claude-single-sonnet}{sn-none}{1}
\setres{fresh-agents}{codex-agent-astra}{s-function}{90.6}
\setres{fresh-agents}{codex-agent-astra}{sn-function}{64}
\setres{fresh-agents}{codex-agent-astra}{s-file}{100.0}
\setres{fresh-agents}{codex-agent-astra}{sn-file}{14}
\setres{fresh-agents}{codex-agent-astra}{s-trace}{66.7}
\setres{fresh-agents}{codex-agent-astra}{sn-trace}{3}
\setres{fresh-agents}{codex-agent-astra}{s-none}{82.9}
\setres{fresh-agents}{codex-agent-astra}{sn-none}{117}
\setres{fresh-agents}{codex-agent}{s-function}{89.1}
\setres{fresh-agents}{codex-agent}{sn-function}{64}
\setres{fresh-agents}{codex-agent}{s-file}{100.0}
\setres{fresh-agents}{codex-agent}{sn-file}{14}
\setres{fresh-agents}{codex-agent}{s-trace}{0.0}
\setres{fresh-agents}{codex-agent}{sn-trace}{3}
\setres{fresh-agents}{codex-agent}{s-none}{67.5}
\setres{fresh-agents}{codex-agent}{sn-none}{117}
\setres{fresh-agents}{codex-single-astra}{s-function}{84.4}
\setres{fresh-agents}{codex-single-astra}{sn-function}{64}
\setres{fresh-agents}{codex-single-astra}{s-file}{78.6}
\setres{fresh-agents}{codex-single-astra}{sn-file}{14}
\setres{fresh-agents}{codex-single-astra}{s-trace}{0.0}
\setres{fresh-agents}{codex-single-astra}{sn-trace}{3}
\setres{fresh-agents}{codex-single-astra}{s-none}{39.3}
\setres{fresh-agents}{codex-single-astra}{sn-none}{117}
\setres{fresh-agents}{codex-single}{s-function}{78.1}
\setres{fresh-agents}{codex-single}{sn-function}{64}
\setres{fresh-agents}{codex-single}{s-file}{35.7}
\setres{fresh-agents}{codex-single}{sn-file}{14}
\setres{fresh-agents}{codex-single}{s-trace}{0.0}
\setres{fresh-agents}{codex-single}{sn-trace}{3}
\setres{fresh-agents}{codex-single}{s-none}{17.1}
\setres{fresh-agents}{codex-single}{sn-none}{117}
\setres{fresh-agents}{gemini-single}{s-function}{85.9}
\setres{fresh-agents}{gemini-single}{sn-function}{64}
\setres{fresh-agents}{gemini-single}{s-file}{57.1}
\setres{fresh-agents}{gemini-single}{sn-file}{14}
\setres{fresh-agents}{gemini-single}{s-trace}{0.0}
\setres{fresh-agents}{gemini-single}{sn-trace}{3}
\setres{fresh-agents}{gemini-single}{s-none}{33.3}
\setres{fresh-agents}{gemini-single}{sn-none}{117}
\setres{fresh-agents}{open-agent-qwen3-coder-next-fp8}{s-function}{82.8}
\setres{fresh-agents}{open-agent-qwen3-coder-next-fp8}{sn-function}{64}
\setres{fresh-agents}{open-agent-qwen3-coder-next-fp8}{s-file}{64.3}
\setres{fresh-agents}{open-agent-qwen3-coder-next-fp8}{sn-file}{14}
\setres{fresh-agents}{open-agent-qwen3-coder-next-fp8}{s-trace}{0.0}
\setres{fresh-agents}{open-agent-qwen3-coder-next-fp8}{sn-trace}{3}
\setres{fresh-agents}{open-agent-qwen3-coder-next-fp8}{s-none}{49.6}
\setres{fresh-agents}{open-agent-qwen3-coder-next-fp8}{sn-none}{117}
\setres{fresh-agents}{open-single-gpt-oss-120b}{s-function}{60.9}
\setres{fresh-agents}{open-single-gpt-oss-120b}{sn-function}{64}
\setres{fresh-agents}{open-single-gpt-oss-120b}{s-file}{14.3}
\setres{fresh-agents}{open-single-gpt-oss-120b}{sn-file}{14}
\setres{fresh-agents}{open-single-gpt-oss-120b}{s-trace}{0.0}
\setres{fresh-agents}{open-single-gpt-oss-120b}{sn-trace}{3}
\setres{fresh-agents}{open-single-gpt-oss-120b}{s-none}{6.8}
\setres{fresh-agents}{open-single-gpt-oss-120b}{sn-none}{117}
\setres{fresh-agents}{open-single-qwen3-coder-next-fp8}{s-function}{71.9}
\setres{fresh-agents}{open-single-qwen3-coder-next-fp8}{sn-function}{64}
\setres{fresh-agents}{open-single-qwen3-coder-next-fp8}{s-file}{21.4}
\setres{fresh-agents}{open-single-qwen3-coder-next-fp8}{sn-file}{14}
\setres{fresh-agents}{open-single-qwen3-coder-next-fp8}{s-trace}{0.0}
\setres{fresh-agents}{open-single-qwen3-coder-next-fp8}{sn-trace}{3}
\setres{fresh-agents}{open-single-qwen3-coder-next-fp8}{s-none}{14.5}
\setres{fresh-agents}{open-single-qwen3-coder-next-fp8}{sn-none}{117}
\setres{fresh-agents}{swerank-embed-large}{s-function}{82.8}
\setres{fresh-agents}{swerank-embed-large}{sn-function}{64}
\setres{fresh-agents}{swerank-embed-large}{s-file}{92.9}
\setres{fresh-agents}{swerank-embed-large}{sn-file}{14}
\setres{fresh-agents}{swerank-embed-large}{s-trace}{33.3}
\setres{fresh-agents}{swerank-embed-large}{sn-trace}{3}
\setres{fresh-agents}{swerank-embed-large}{s-none}{61.5}
\setres{fresh-agents}{swerank-embed-large}{sn-none}{117}
\setres{fresh-agents}{swerank-embed-small}{s-function}{71.9}
\setres{fresh-agents}{swerank-embed-small}{sn-function}{64}
\setres{fresh-agents}{swerank-embed-small}{s-file}{64.3}
\setres{fresh-agents}{swerank-embed-small}{sn-file}{14}
\setres{fresh-agents}{swerank-embed-small}{s-trace}{0.0}
\setres{fresh-agents}{swerank-embed-small}{sn-trace}{3}
\setres{fresh-agents}{swerank-embed-small}{s-none}{59.0}
\setres{fresh-agents}{swerank-embed-small}{sn-none}{117}
\setres{fresh-agents}{website-gpu}{s-function}{95.3}
\setres{fresh-agents}{website-gpu}{sn-function}{64}
\setres{fresh-agents}{website-gpu}{s-file}{92.9}
\setres{fresh-agents}{website-gpu}{sn-file}{14}
\setres{fresh-agents}{website-gpu}{s-trace}{0.0}
\setres{fresh-agents}{website-gpu}{sn-trace}{3}
\setres{fresh-agents}{website-gpu}{s-none}{62.4}
\setres{fresh-agents}{website-gpu}{sn-none}{117}
\setres{fresh-agents}{website}{s-function}{95.3}
\setres{fresh-agents}{website}{sn-function}{64}
\setres{fresh-agents}{website}{s-file}{92.9}
\setres{fresh-agents}{website}{sn-file}{14}
\setres{fresh-agents}{website}{s-trace}{0.0}
\setres{fresh-agents}{website}{sn-trace}{3}
\setres{fresh-agents}{website}{s-none}{61.5}
\setres{fresh-agents}{website}{sn-none}{117}
\setres{fresh-agents}{claude-agent}{esc-escalated}{62.6}
\setres{fresh-agents}{claude-agent}{esc-hit@10}{91.4}
\setres{fresh-agents}{claude-agent}{esc-random_hit@10}{84.7}
\setres{fresh-agents}{claude-single-sonnet}{esc-escalated}{100.0}
\setres{fresh-agents}{claude-single-sonnet}{esc-hit@10}{100.0}
\setres{fresh-agents}{claude-single-sonnet}{esc-random_hit@10}{100.0}
\setres{fresh-agents}{codex-agent-astra}{esc-escalated}{62.6}
\setres{fresh-agents}{codex-agent-astra}{esc-hit@10}{86.9}
\setres{fresh-agents}{codex-agent-astra}{esc-random_hit@10}{81.8}
\setres{fresh-agents}{codex-agent}{esc-escalated}{62.6}
\setres{fresh-agents}{codex-agent}{esc-hit@10}{78.3}
\setres{fresh-agents}{codex-agent}{esc-random_hit@10}{75.2}
\setres{fresh-agents}{codex-single-astra}{esc-escalated}{62.6}
\setres{fresh-agents}{codex-single-astra}{esc-hit@10}{62.6}
\setres{fresh-agents}{codex-single-astra}{esc-random_hit@10}{62.9}
\setres{fresh-agents}{codex-single}{esc-escalated}{62.6}
\setres{fresh-agents}{codex-single}{esc-hit@10}{49.0}
\setres{fresh-agents}{codex-single}{esc-random_hit@10}{51.5}
\setres{fresh-agents}{gemini-single}{esc-escalated}{62.6}
\setres{fresh-agents}{gemini-single}{esc-hit@10}{57.6}
\setres{fresh-agents}{gemini-single}{esc-random_hit@10}{60.0}
\setres{fresh-agents}{open-agent-qwen3-coder-next-fp8}{esc-escalated}{62.6}
\setres{fresh-agents}{open-agent-qwen3-coder-next-fp8}{esc-hit@10}{65.7}
\setres{fresh-agents}{open-agent-qwen3-coder-next-fp8}{esc-random_hit@10}{65.7}
\setres{fresh-agents}{open-single-gpt-oss-120b}{esc-escalated}{62.6}
\setres{fresh-agents}{open-single-gpt-oss-120b}{esc-hit@10}{40.4}
\setres{fresh-agents}{open-single-gpt-oss-120b}{esc-random_hit@10}{43.2}
\setres{fresh-agents}{open-single-qwen3-coder-next-fp8}{esc-escalated}{62.6}
\setres{fresh-agents}{open-single-qwen3-coder-next-fp8}{esc-hit@10}{47.5}
\setres{fresh-agents}{open-single-qwen3-coder-next-fp8}{esc-random_hit@10}{48.6}
\setres{fresh-agents}{swerank-embed-large}{esc-escalated}{62.6}
\setres{fresh-agents}{swerank-embed-large}{esc-hit@10}{75.8}
\setres{fresh-agents}{swerank-embed-large}{esc-random_hit@10}{71.7}
\setres{fresh-agents}{swerank-embed-small}{esc-escalated}{62.6}
\setres{fresh-agents}{swerank-embed-small}{esc-hit@10}{70.2}
\setres{fresh-agents}{swerank-embed-small}{esc-random_hit@10}{67.0}
\setres{fresh-agents}{website}{esc-escalated}{62.6}
\setres{fresh-agents}{website}{esc-hit@10}{73.7}
\setres{fresh-agents}{website}{esc-random_hit@10}{73.9}
\setres{fresh}{allfather-fast}{hit@1}{62.0}
\setres{fresh}{allfather-fast}{hit@1-ci}{50.4--71.5}
\setres{fresh}{allfather-fast}{hit@5}{83.0}
\setres{fresh}{allfather-fast}{hit@5-ci}{74.5--90.6}
\setres{fresh}{allfather-fast}{hit@10}{88.0}
\setres{fresh}{allfather-fast}{hit@10-ci}{80.0--94.3}
\setres{fresh}{allfather-fast}{all@10}{54.0}
\setres{fresh}{allfather-fast}{all@10-ci}{47.0--59.8}
\setres{fresh}{allfather-fast}{mrr}{0.708}
\setres{fresh}{allfather-fast}{n}{100}
\setres{fresh}{allfather-fast}{cost}{\$0}
\setres{fresh}{allfather-fast}{sec}{0.0}
\setres{fresh}{claude-single}{hit@1}{57.1}
\setres{fresh}{claude-single}{hit@1-ci}{47.1--67.2}
\setres{fresh}{claude-single}{hit@5}{71.4}
\setres{fresh}{claude-single}{hit@5-ci}{59.8--84.4}
\setres{fresh}{claude-single}{hit@10}{74.0}
\setres{fresh}{claude-single}{hit@10-ci}{62.5--85.5}
\setres{fresh}{claude-single}{all@10}{42.9}
\setres{fresh}{claude-single}{all@10-ci}{38.2--48.8}
\setres{fresh}{claude-single}{mrr}{0.635}
\setres{fresh}{claude-single}{n}{77}
\setres{fresh}{claude-single}{cost}{\$1.16}
\setres{fresh}{claude-single}{sec}{8.4}
\setres{fresh}{codex-single-astra}{hit@1}{39.7}
\setres{fresh}{codex-single-astra}{hit@1-ci}{34.1--45.1}
\setres{fresh}{codex-single-astra}{hit@5}{50.0}
\setres{fresh}{codex-single-astra}{hit@5-ci}{43.7--55.9}
\setres{fresh}{codex-single-astra}{hit@10}{50.0}
\setres{fresh}{codex-single-astra}{hit@10-ci}{43.7--55.9}
\setres{fresh}{codex-single-astra}{all@10}{20.9}
\setres{fresh}{codex-single-astra}{all@10-ci}{16.7--25.1}
\setres{fresh}{codex-single-astra}{mrr}{0.441}
\setres{fresh}{codex-single-astra}{n}{580}
\setres{fresh}{codex-single-astra}{cost}{\$0.49}
\setres{fresh}{codex-single-astra}{sec}{15.2}
\setres{fresh}{codex-single}{hit@1}{28.1}
\setres{fresh}{codex-single}{hit@1-ci}{22.6--33.7}
\setres{fresh}{codex-single}{hit@5}{32.8}
\setres{fresh}{codex-single}{hit@5-ci}{27.3--38.2}
\setres{fresh}{codex-single}{hit@10}{32.8}
\setres{fresh}{codex-single}{hit@10-ci}{27.3--38.2}
\setres{fresh}{codex-single}{all@10}{15.2}
\setres{fresh}{codex-single}{all@10-ci}{11.6--18.8}
\setres{fresh}{codex-single}{mrr}{0.303}
\setres{fresh}{codex-single}{n}{580}
\setres{fresh}{codex-single}{cost}{\$0.005}
\setres{fresh}{codex-single}{sec}{9.9}
\setres{fresh}{swerank-embed-large}{hit@1}{43.3}
\setres{fresh}{swerank-embed-large}{hit@1-ci}{38.1--48.4}
\setres{fresh}{swerank-embed-large}{hit@5}{65.5}
\setres{fresh}{swerank-embed-large}{hit@5-ci}{59.2--71.4}
\setres{fresh}{swerank-embed-large}{hit@10}{70.9}
\setres{fresh}{swerank-embed-large}{hit@10-ci}{64.7--76.6}
\setres{fresh}{swerank-embed-large}{all@10}{36.9}
\setres{fresh}{swerank-embed-large}{all@10-ci}{31.6--41.8}
\setres{fresh}{swerank-embed-large}{mrr}{0.529}
\setres{fresh}{swerank-embed-large}{n}{580}
\setres{fresh}{swerank-embed-large}{cost}{--}
\setres{fresh}{swerank-embed-large}{sec}{0.0}
\setres{fresh}{swerank-embed-small}{hit@1}{38.3}
\setres{fresh}{swerank-embed-small}{hit@1-ci}{32.8--43.5}
\setres{fresh}{swerank-embed-small}{hit@5}{57.4}
\setres{fresh}{swerank-embed-small}{hit@5-ci}{50.8--63.7}
\setres{fresh}{swerank-embed-small}{hit@10}{65.3}
\setres{fresh}{swerank-embed-small}{hit@10-ci}{58.7--71.7}
\setres{fresh}{swerank-embed-small}{all@10}{34.0}
\setres{fresh}{swerank-embed-small}{all@10-ci}{28.9--38.9}
\setres{fresh}{swerank-embed-small}{mrr}{0.465}
\setres{fresh}{swerank-embed-small}{n}{580}
\setres{fresh}{swerank-embed-small}{cost}{--}
\setres{fresh}{swerank-embed-small}{sec}{0.0}
\setres{fresh}{website-gpu}{hit@1}{47.1}
\setres{fresh}{website-gpu}{hit@1-ci}{42.3--51.6}
\setres{fresh}{website-gpu}{hit@5}{70.5}
\setres{fresh}{website-gpu}{hit@5-ci}{65.7--75.2}
\setres{fresh}{website-gpu}{hit@10}{76.0}
\setres{fresh}{website-gpu}{hit@10-ci}{71.4--80.4}
\setres{fresh}{website-gpu}{all@10}{42.9}
\setres{fresh}{website-gpu}{all@10-ci}{37.7--47.9}
\setres{fresh}{website-gpu}{mrr}{0.568}
\setres{fresh}{website-gpu}{n}{580}
\setres{fresh}{website-gpu}{cost}{\$0.001}
\setres{fresh}{website-gpu}{sec}{7.3}
\setres{fresh}{website}{hit@1}{42.4}
\setres{fresh}{website}{hit@1-ci}{37.3--47.4}
\setres{fresh}{website}{hit@5}{65.0}
\setres{fresh}{website}{hit@5-ci}{59.8--69.7}
\setres{fresh}{website}{hit@10}{71.6}
\setres{fresh}{website}{hit@10-ci}{66.5--76.2}
\setres{fresh}{website}{all@10}{37.8}
\setres{fresh}{website}{all@10-ci}{33.0--42.3}
\setres{fresh}{website}{mrr}{0.521}
\setres{fresh}{website}{n}{580}
\setres{fresh}{website}{cost}{\$0}
\setres{fresh}{website}{sec}{4.6}
\setres{fresh}{allfather-fast}{dhit@1}{+3.0}
\setres{fresh}{allfather-fast}{wlhit@1}{9/6}
\setres{fresh}{allfather-fast}{phit@1}{=1.000}
\setres{fresh}{allfather-fast}{dhit@10}{+5.0}
\setres{fresh}{allfather-fast}{wlhit@10}{6/1}
\setres{fresh}{allfather-fast}{phit@10}{=1.000}
\setres{fresh}{claude-single}{dhit@1}{+1.3}
\setres{fresh}{claude-single}{wlhit@1}{12/11}
\setres{fresh}{claude-single}{phit@1}{=1.000}
\setres{fresh}{claude-single}{dhit@10}{-6.5}
\setres{fresh}{claude-single}{wlhit@10}{5/10}
\setres{fresh}{claude-single}{phit@10}{=1.000}
\setres{fresh}{codex-single-astra}{dhit@1}{-7.4}
\setres{fresh}{codex-single-astra}{wlhit@1}{74/117}
\setres{fresh}{codex-single-astra}{phit@1}{=0.037}
\setres{fresh}{codex-single-astra}{dhit@10}{-26.0}
\setres{fresh}{codex-single-astra}{wlhit@10}{31/182}
\setres{fresh}{codex-single-astra}{phit@10}{<0.001}
\setres{fresh}{codex-single}{dhit@1}{-19.0}
\setres{fresh}{codex-single}{wlhit@1}{35/145}
\setres{fresh}{codex-single}{phit@1}{<0.001}
\setres{fresh}{codex-single}{dhit@10}{-43.3}
\setres{fresh}{codex-single}{wlhit@10}{12/263}
\setres{fresh}{codex-single}{phit@10}{<0.001}
\setres{fresh}{swerank-embed-large}{dhit@1}{-3.8}
\setres{fresh}{swerank-embed-large}{wlhit@1}{67/89}
\setres{fresh}{swerank-embed-large}{phit@1}{=0.924}
\setres{fresh}{swerank-embed-large}{dhit@10}{-5.2}
\setres{fresh}{swerank-embed-large}{wlhit@10}{40/70}
\setres{fresh}{swerank-embed-large}{phit@10}{=0.076}
\setres{fresh}{swerank-embed-small}{dhit@1}{-8.8}
\setres{fresh}{swerank-embed-small}{wlhit@1}{53/104}
\setres{fresh}{swerank-embed-small}{phit@1}{=0.002}
\setres{fresh}{swerank-embed-small}{dhit@10}{-10.7}
\setres{fresh}{swerank-embed-small}{wlhit@10}{34/96}
\setres{fresh}{swerank-embed-small}{phit@10}{<0.001}
\setres{fresh}{website}{dhit@1}{-4.7}
\setres{fresh}{website}{wlhit@1}{59/86}
\setres{fresh}{website}{phit@1}{=0.336}
\setres{fresh}{website}{dhit@10}{-4.5}
\setres{fresh}{website}{wlhit@10}{11/37}
\setres{fresh}{website}{phit@10}{=0.004}
\setres{fresh}{allfather-fast}{browser-dhit@1}{+9.0}
\setres{fresh}{allfather-fast}{browser-wlhit@1}{18/9}
\setres{fresh}{allfather-fast}{browser-phit@1}{=1.000}
\setres{fresh}{allfather-fast}{browser-dhit@10}{+6.0}
\setres{fresh}{allfather-fast}{browser-wlhit@10}{7/1}
\setres{fresh}{allfather-fast}{browser-phit@10}{=1.000}
\setres{fresh}{claude-single}{browser-dhit@1}{+1.3}
\setres{fresh}{claude-single}{browser-wlhit@1}{13/12}
\setres{fresh}{claude-single}{browser-phit@1}{=1.000}
\setres{fresh}{claude-single}{browser-dhit@10}{+1.3}
\setres{fresh}{claude-single}{browser-wlhit@10}{8/7}
\setres{fresh}{claude-single}{browser-phit@10}{=1.000}
\setres{fresh}{codex-single-astra}{browser-dhit@1}{-2.8}
\setres{fresh}{codex-single-astra}{browser-wlhit@1}{82/98}
\setres{fresh}{codex-single-astra}{browser-phit@1}{=1.000}
\setres{fresh}{codex-single-astra}{browser-dhit@10}{-21.6}
\setres{fresh}{codex-single-astra}{browser-wlhit@10}{37/162}
\setres{fresh}{codex-single-astra}{browser-phit@10}{<0.001}
\setres{fresh}{codex-single}{browser-dhit@1}{-14.3}
\setres{fresh}{codex-single}{browser-wlhit@1}{50/133}
\setres{fresh}{codex-single}{browser-phit@1}{<0.001}
\setres{fresh}{codex-single}{browser-dhit@10}{-38.8}
\setres{fresh}{codex-single}{browser-wlhit@10}{15/240}
\setres{fresh}{codex-single}{browser-phit@10}{<0.001}
\setres{fresh}{swerank-embed-large}{browser-dhit@1}{+0.9}
\setres{fresh}{swerank-embed-large}{browser-wlhit@1}{78/73}
\setres{fresh}{swerank-embed-large}{browser-phit@1}{=1.000}
\setres{fresh}{swerank-embed-large}{browser-dhit@10}{-0.7}
\setres{fresh}{swerank-embed-large}{browser-wlhit@10}{59/63}
\setres{fresh}{swerank-embed-large}{browser-phit@10}{=1.000}
\setres{fresh}{swerank-embed-small}{browser-dhit@1}{-4.1}
\setres{fresh}{swerank-embed-small}{browser-wlhit@1}{63/87}
\setres{fresh}{swerank-embed-small}{browser-phit@1}{=0.960}
\setres{fresh}{swerank-embed-small}{browser-dhit@10}{-6.2}
\setres{fresh}{swerank-embed-small}{browser-wlhit@10}{45/81}
\setres{fresh}{swerank-embed-small}{browser-phit@10}{=0.034}
\setres{fresh}{website-gpu}{browser-dhit@1}{+4.7}
\setres{fresh}{website-gpu}{browser-wlhit@1}{86/59}
\setres{fresh}{website-gpu}{browser-phit@1}{=0.549}
\setres{fresh}{website-gpu}{browser-dhit@10}{+4.5}
\setres{fresh}{website-gpu}{browser-wlhit@10}{37/11}
\setres{fresh}{website-gpu}{browser-phit@10}{=0.005}
\setres{fresh}{allfather-fast}{s-function}{100.0}
\setres{fresh}{allfather-fast}{sn-function}{34}
\setres{fresh}{allfather-fast}{s-file}{100.0}
\setres{fresh}{allfather-fast}{sn-file}{12}
\setres{fresh}{allfather-fast}{s-trace}{0.0}
\setres{fresh}{allfather-fast}{sn-trace}{1}
\setres{fresh}{allfather-fast}{s-none}{79.2}
\setres{fresh}{allfather-fast}{sn-none}{53}
\setres{fresh}{claude-single}{s-function}{97.0}
\setres{fresh}{claude-single}{sn-function}{33}
\setres{fresh}{claude-single}{s-file}{54.5}
\setres{fresh}{claude-single}{sn-file}{11}
\setres{fresh}{claude-single}{s-trace}{100.0}
\setres{fresh}{claude-single}{sn-trace}{1}
\setres{fresh}{claude-single}{s-none}{56.2}
\setres{fresh}{claude-single}{sn-none}{32}
\setres{fresh}{codex-single-astra}{s-function}{77.4}
\setres{fresh}{codex-single-astra}{sn-function}{186}
\setres{fresh}{codex-single-astra}{s-file}{56.0}
\setres{fresh}{codex-single-astra}{sn-file}{50}
\setres{fresh}{codex-single-astra}{s-trace}{0.0}
\setres{fresh}{codex-single-astra}{sn-trace}{5}
\setres{fresh}{codex-single-astra}{s-none}{34.8}
\setres{fresh}{codex-single-astra}{sn-none}{339}
\setres{fresh}{codex-single}{s-function}{69.9}
\setres{fresh}{codex-single}{sn-function}{186}
\setres{fresh}{codex-single}{s-file}{28.0}
\setres{fresh}{codex-single}{sn-file}{50}
\setres{fresh}{codex-single}{s-trace}{0.0}
\setres{fresh}{codex-single}{sn-trace}{5}
\setres{fresh}{codex-single}{s-none}{13.6}
\setres{fresh}{codex-single}{sn-none}{339}
\setres{fresh}{swerank-embed-large}{s-function}{84.4}
\setres{fresh}{swerank-embed-large}{sn-function}{186}
\setres{fresh}{swerank-embed-large}{s-file}{84.0}
\setres{fresh}{swerank-embed-large}{sn-file}{50}
\setres{fresh}{swerank-embed-large}{s-trace}{40.0}
\setres{fresh}{swerank-embed-large}{sn-trace}{5}
\setres{fresh}{swerank-embed-large}{s-none}{62.0}
\setres{fresh}{swerank-embed-large}{sn-none}{339}
\setres{fresh}{swerank-embed-small}{s-function}{81.7}
\setres{fresh}{swerank-embed-small}{sn-function}{186}
\setres{fresh}{swerank-embed-small}{s-file}{72.0}
\setres{fresh}{swerank-embed-small}{sn-file}{50}
\setres{fresh}{swerank-embed-small}{s-trace}{20.0}
\setres{fresh}{swerank-embed-small}{sn-trace}{5}
\setres{fresh}{swerank-embed-small}{s-none}{56.0}
\setres{fresh}{swerank-embed-small}{sn-none}{339}
\setres{fresh}{website-gpu}{s-function}{96.2}
\setres{fresh}{website-gpu}{sn-function}{186}
\setres{fresh}{website-gpu}{s-file}{98.0}
\setres{fresh}{website-gpu}{sn-file}{50}
\setres{fresh}{website-gpu}{s-trace}{0.0}
\setres{fresh}{website-gpu}{sn-trace}{5}
\setres{fresh}{website-gpu}{s-none}{62.8}
\setres{fresh}{website-gpu}{sn-none}{339}
\setres{fresh}{website}{s-function}{94.6}
\setres{fresh}{website}{sn-function}{186}
\setres{fresh}{website}{s-file}{98.0}
\setres{fresh}{website}{sn-file}{50}
\setres{fresh}{website}{s-trace}{0.0}
\setres{fresh}{website}{sn-trace}{5}
\setres{fresh}{website}{s-none}{56.0}
\setres{fresh}{website}{sn-none}{339}
\setres{fresh}{allfather-fast}{esc-escalated}{54.0}
\setres{fresh}{allfather-fast}{esc-hit@10}{88.0}
\setres{fresh}{allfather-fast}{esc-random_hit@10}{85.7}
\setres{fresh}{claude-single}{esc-escalated}{45.5}
\setres{fresh}{claude-single}{esc-hit@10}{80.5}
\setres{fresh}{claude-single}{esc-random_hit@10}{77.6}
\setres{fresh}{codex-single-astra}{esc-escalated}{62.8}
\setres{fresh}{codex-single-astra}{esc-hit@10}{59.0}
\setres{fresh}{codex-single-astra}{esc-random_hit@10}{59.7}
\setres{fresh}{codex-single}{esc-escalated}{62.8}
\setres{fresh}{codex-single}{esc-hit@10}{46.0}
\setres{fresh}{codex-single}{esc-random_hit@10}{48.9}
\setres{fresh}{swerank-embed-large}{esc-escalated}{62.8}
\setres{fresh}{swerank-embed-large}{esc-hit@10}{75.9}
\setres{fresh}{swerank-embed-large}{esc-random_hit@10}{72.8}
\setres{fresh}{swerank-embed-small}{esc-escalated}{62.8}
\setres{fresh}{swerank-embed-small}{esc-hit@10}{71.0}
\setres{fresh}{swerank-embed-small}{esc-random_hit@10}{69.3}
\setres{fresh}{website}{esc-escalated}{62.8}
\setres{fresh}{website}{esc-hit@10}{72.2}
\setres{fresh}{website}{esc-random_hit@10}{73.2}


\newif\ifanon
\anontrue
\ifanon
  \newcommand{\tool}{\textsc{Loc}\xspace}
  \newcommand{\artifact}{\url{https://anonymous.4open.science/r/fresh-2026}}
\else
  \newcommand{\tool}{Allfather\xspace}
  \newcommand{\artifact}{\url{https://github.com/Lazeeez/Allfather}}
\fi
\newcommand{\toolgpu}{\tool{}+GPU\xspace}
\newcommand{\fresh}{\textsc{Fresh-2026}\xspace}
\newcommand{\FF}{fresh-full-agents}
\newcommand{\FFA}{fresh-full}
\newcommand{\SH}{fresh-agents}
\newcommand{\SHA}{fresh}
\newcommand{\RQ}[1]{\textbf{RQ#1}}
\newcommand{\pending}{\textbf{\textcolor{red}{[pending]}}}

\newtcolorbox{finding}[1]{enhanced, breakable, colback=gray!6, colframe=gray!60, boxrule=0.5pt, arc=1pt,
  left=4pt, right=4pt, top=3pt, bottom=3pt, fonttitle=\bfseries\small, title={#1}, fontupper=\small}

\setlength{\emergencystretch}{1.5em}

\begin{document}

\title{Do You Need a Frontier Agent to Find the Bug? A Pre-Registered Comparison of Bug Localizers on Post-Cutoff Bugs}

\author{Lajat Manekar}
\affiliation{\institution{Georgia Institute of Technology}\city{Atlanta}\state{Georgia}\country{USA}}
\email{lajatm9@gmail.com}

\begin{abstract}
Coding agents built on frontier language models now localize bugs by reading a repository with tools, and they report strong results on SWE-bench. Those results are hard to interpret for two reasons: the benchmark's bugs predate the models, which may have seen the fixes during training, and the comparisons rarely include cost, privacy, or the tools a developer can use for free. I present \fresh, a benchmark of \NBugs{} bugs reported between 1 August and 2 November 2026 in \NProjects{} open-source projects written in Go, Java, JavaScript, Python, Rust, and TypeScript. Every bug postdates the training cutoff of every model studied, was vetted for bot- and AI-written issues and fixes, and was sealed by hash before any system ran. On it, I compare a free bug localizer that runs entirely in the browser, with an optional small GPU reranker, against a published retrieval system, two open-weight models, free-tier and paid frontier models asked once, and frontier coding agents, under a protocol fixed in advance. With full bug reports, the free localizer reaches \res{\FFA}{website-gpu}{hit@10}\% function-level hit@10 on all \NBugs{} bugs, more than any other system that ran on the whole set. On a \NAgent{}-bug subset, it ties the best free coding agent at hit@10 and outperforms both open-weight models, as agents and as single calls, and every single-call model; the agent, however, ranks the faulty function first far more often. On shortened reports, paid frontier agents are more accurate still, and a confidence gate fixed before the study recovers most of their advantage while escalating only part of the bugs. Finally, frontier models asked once fall far behind the same models used as agents on \fresh, a gap that is absent on SWE-bench Verified and is consistent with memorization inflating older benchmarks.
\end{abstract}

\begin{CCSXML}
<ccs2012>
<concept><concept_id>10011007.10011074.10011099.10011102.10011103</concept_id><concept_desc>Software and its engineering~Software testing and debugging</concept_desc><concept_significance>500</concept_significance></concept>
<concept><concept_id>10011007.10011006.10011073</concept_id><concept_desc>Software and its engineering~Software maintenance tools</concept_desc><concept_significance>300</concept_significance></concept>
</ccs2012>
\end{CCSXML}
\ccsdesc[500]{Software and its engineering~Software testing and debugging}
\ccsdesc[300]{Software and its engineering~Software maintenance tools}

\keywords{bug localization, fault localization, large language models, coding agents, benchmark contamination, empirical study}

\maketitle

%% ---------------------------------------------------------------------------
\section{Introduction}
\label{sec:intro}

Before a developer can fix a reported bug, they must find the code responsible for it. Bug localization automates this step: given a bug report and a repository, it ranks the repository's files or functions by how likely they are to need a change~\cite{zhou2012buglocator,saha2013bluir,ye2014learning}. For two decades the field relied on information retrieval and learning to rank, evaluated mostly on a handful of Java projects~\cite{lee2018bench4bl}. In the last two years, large language models have changed both the methods and the evaluation. Coding agents now localize bugs the way a developer would, by listing, searching, and reading files~\cite{yang2024sweagent,xia2024agentless,chen2025locagent}, and they are evaluated on SWE-bench~\cite{jimenez2024swebench}, a set of GitHub issues from twelve popular Python projects.

A developer who wants to adopt one of these tools faces three open questions. The first is whether the reported accuracy holds on new bugs. SWE-bench's issues were filed between 2012 and 2023, before any current model was trained, and frontier models have read the same public repositories, issues, and fixes. When a model has seen a fix, a benchmark measures recall rather than localization. In a ten-bug pilot on SWE-bench Verified, Claude Opus~5.5 given only the bug report and the list of file paths, with no way to read any code, localized as many bugs as the same model running as an agent that could read every file. The second question is cost. Agents spend tokens on every bug, and published comparisons rarely report what a correct answer costs, or compare against tools that cost nothing. The third question is privacy: every model-based tool sends the repository and the report to a third party, which many organizations cannot do for proprietary code.

This paper answers these questions with a pre-registered empirical study. I built \fresh, \NBugs{} bugs from \NProjects{} open-source projects in six languages, all reported after the published training cutoff of every model in the study and sealed before any system ran. On it I compare sixteen system configurations, from a localizer that runs in the user's browser to frontier coding agents, on the same bugs, with the same answer format, and with the cost of every answer measured. The protocol, including the systems, metrics, statistical tests, and decision thresholds, was written before the test set was scored; every later change is a dated amendment written before the results it affects were read. I study four research questions:

\begin{description}
\item[\RQ1 (accuracy)] On bugs no model can have seen, how accurate is each kind of system relative to a free, browser-based localizer?
\item[\RQ2 (cost)] Which systems are not outperformed by any other system on both accuracy and cost?
\item[\RQ3 (escalation)] Can a confidence gate that sends only uncertain bugs to an agent approach the agent's accuracy at a fraction of its cost?
\item[\RQ4 (where the gap lies)] How do the systems differ by programming language and by the information a bug report contains?
\end{description}

The free localizer studied here, \tool, splits a repository into functions, indexes them with a code-aware lexical index and static embeddings, and ranks them with a learned ranker; all of it runs in a web browser, and no code leaves the device. An optional variant, \toolgpu, re-reads the top 50 candidates with a 0.6B-parameter cross-encoder on a GPU. The main findings are as follows.

\begin{itemize}
\item \textbf{Among free systems, \tool is tied for most accurate at hit@10.} On all \NBugs{} full bug reports, \toolgpu places a changed function in its top ten for \res{\FFA}{website-gpu}{hit@10}\% of bugs, ahead of SweRankEmbed-Large (\res{\FFA}{swerank-embed-large}{hit@10}\%), the strongest published retriever. On the \NAgent{}-bug subset, it ties the best free coding agent, gpt-6-luna in Codex (\res{\FF}{website-gpu}{hit@10}\% against \res{\FF}{codex-agent}{hit@10}\%), and outperforms both open-weight models and every single-call model.
\item \textbf{Agents rank better.} The free agent puts a changed function first for \res{\FF}{codex-agent}{hit@1}\% of bugs, against \res{\FF}{website-gpu}{hit@1}\% for \toolgpu.
\item \textbf{Paid frontier agents are the most accurate systems, and a gate can call them sparingly.} On shortened reports, Claude Opus~5.5 in Claude Code reaches \ClaudeAgentShort\% hit@10. Escalating to it only the \CascadeClaudeEscalated\% of bugs \tool is least confident about yields \CascadeClaude\%, against \CascadeClaudeRandom\% when the same number of bugs is escalated at random.
\item \textbf{Single-call accuracy collapses on fresh bugs.} gpt-6-astra asked once reaches \AstraSingleShort\% on the bugs where its agent reaches \AstraAgentShort\%, although on SWE-bench Verified the single call matched the agent.
\end{itemize}

\noindent In summary, this paper makes the following contributions:
\begin{itemize}
\item \fresh, a six-language bug-localization benchmark of post-cutoff bugs with a documented vetting and sealing procedure (Section~\ref{sec:benchmark});
\item a pre-registered comparison of sixteen configurations, including frontier agents used as developers use them, with exact paired tests, project-clustered intervals, and per-answer cost (Sections~\ref{sec:design} and~\ref{sec:results});
\item evidence that a free, private-by-default localizer matches the best free coding agent at hit@10 on unseen bugs, together with an account of where it falls short (Section~\ref{sec:rq1});
\item a measured design for escalating only uncertain bugs to a paid agent (Section~\ref{sec:rq3});
\item all data, run records, prompts, and analysis scripts, from which every number in this paper is regenerated (Data Availability).
\end{itemize}

%% ---------------------------------------------------------------------------
\section{Background}
\label{sec:background}

\subsection{Bug Localization}

Bug localization takes a bug report, usually a title and a free-text description, and a repository at the commit where the bug was reported, and returns a ranked list of code locations. The ground truth is the set of locations the developers' fix changed. Localization can be evaluated at file, function, or line granularity; function level is a useful middle ground, because a function is small enough to inspect and is the unit most fixes change~\cite{chen2025locagent}. A localizer is useful to a developer only if the faulty location appears near the top of the list: practitioners report that they would inspect at most the first five to ten results~\cite{kochhar2016practitioners}, and a correct answer far down the list may not help at all~\cite{parnin2011automated}.

\subsection{Three Families of Localizers}

\textbf{Retrieval and learning to rank.} Classic localizers treat the report as a query and code as documents, scoring them with vector-space or BM25 models, sometimes with the history of similar bugs~\cite{zhou2012buglocator,saha2013bluir}, and combine many such signals with learning to rank~\cite{ye2014learning}. Neural retrievers embed reports and code in one space, and cross-encoders rerank a short list by reading the report and the code together~\cite{nogueira2019passage}. SweRank trains both kinds of model for issue localization on pairs mined from GitHub~\cite{reddy2025swerank}.

\textbf{Single-call language models.} A language model can be given the report and a description of the repository, such as its file list, and asked for the faulty locations in one response. This is cheap, but the model cannot read code it is not shown, so its answer depends on what it can infer from names and on what it remembers.

\textbf{Agents.} A coding agent gives a model tools to list, search, and read files and lets it decide what to inspect, over many turns~\cite{yang2024sweagent}. Agentless localizes hierarchically, from files to functions, without an open-ended loop~\cite{xia2024agentless}; LocAgent searches a graph of the code~\cite{chen2025locagent}. Commercial agents such as Claude Code and Codex make this workflow available to every developer.

\subsection{Contamination}

A benchmark drawn from public repositories can leak into a model's training data. When it does, a model may reproduce a known fix without localizing anything, and the benchmark overstates its ability on new bugs. The usual remedy is a benchmark built only from data created after the models' training cutoffs~\cite{jain2024livecodebench}. SWE-bench's issues predate every current frontier model, so its results cannot separate localization from recall.

%% ---------------------------------------------------------------------------
\section{The \fresh Benchmark}
\label{sec:benchmark}

\subsection{Construction}

\textbf{Projects.} I selected public GitHub projects in the six languages, taking the most-starred active projects with issues enabled. I excluded every project used in any of \tool's training, development, or earlier test sets, comparing repositories by GitHub's numeric identifier so that renamed or forked copies are also caught.

\textbf{Bugs.} A bug enters the candidate pool if its issue was created on or after 1 August 2026, after the latest published training cutoff of the compared models, and was closed by exactly one merged pull request that was opened after the issue. The ground truth is the set of functions that the merged diff changes, excluding tests, at the pull request's base commit. GitHub serves only an issue's current text, so issues edited after their fix's pull request was opened are excluded, as is any issue whose fix changes no function. Of \NCandidates{} candidate issues, \NUsable{} met these criteria.

\textbf{Sampling.} I drew up to 100 bugs per language at random with a fixed seed, with at most 12 bugs from any one project so that no project dominates a language. Java had fewer eligible bugs than the other languages and contributes 80. Table~\ref{tab:dataset} summarizes the result.

\begin{table}[t]
\centering
\small
\caption{fresh-2026: bugs reported from 1 August to 2 November 2026.}
\label{tab:dataset}
\resizebox{\columnwidth}{!}{\begin{tabular}{lrrrrr}
\toprule
Language & Bugs & Projects & Agent subset & Funcs/bug & Report chars (full) \\
\midrule
Python & 100 & 16 & 34 & 2 & 2,826 \\
JavaScript & 100 & 20 & 34 & 2 & 2,748 \\
TypeScript & 100 & 20 & 34 & 2 & 2,096 \\
Go & 100 & 16 & 34 & 3 & 3,326 \\
Java & 80 & 12 & 28 & 2 & 1,958 \\
Rust & 100 & 15 & 34 & 1 & 2,086 \\
\midrule
All & 580 & 74 & 198 & 2 & 2,554 \\
\bottomrule
\end{tabular}}
\end{table}


\textbf{Vetting.} A bug written by an AI model, or fixed by one, would favor the models under study. Before sealing, I fetched each drawn bug's issue and pull request again and excluded the bug if either was opened by a bot or a known coding-agent account, or if either stated that it was AI-written (for example, ``Generated with Claude Code''), drawing the next random bug in its place. Of \NVetted{} drawn bugs, \NVetExcluded{} were excluded: \NVetAI{} because the issue or pull request said it was AI-written, and \NVetBot{} because a bot or agent account opened it. AI assistance that leaves no trace cannot be detected and remains a threat to validity (Section~\ref{sec:threats}).

\textbf{Sealing.} I committed the benchmark and the SHA-256 hash of its ground truth (beginning \texttt{\GoldHash}) on 3 November 2026, before any system ran on it, and every later run is recorded in a log in the artifact. The \NAgent{}-bug agent subset, a random sample stratified by language, was fixed at the same time.

\subsection{Report Conditions}

The protocol first gave every system the issue's title and the first 1,500 characters of its body, a common way to bound prompt length. \NLonger{} of the \NBugs{} bodies are longer than that. Because a truncated report may hide the very information a localizer needs, I added a second condition, the full title and body, by an amendment written before the full texts were fetched or any result read. The full report is the primary condition in this paper; the shortened condition is reported as a robustness check (Section~\ref{sec:robustness}).

\subsection{Hint Strata}

What a report says strongly affects how hard a bug is to localize. I assign each bug to the first of four strata that applies: the report names a function the fix changed (\emph{function}); otherwise, it names a changed file (\emph{file}); otherwise, it contains a stack trace (\emph{trace}); otherwise, \emph{none}. A parser assigns the strata. To check it, I labeled 50 randomly drawn bugs by hand without seeing the parser's labels, and I report Cohen's kappa~\cite{cohen1960kappa}; the protocol treats the strata analysis as indicative only if kappa is below 0.6.

%% ---------------------------------------------------------------------------
\section{Study Design}
\label{sec:design}

\subsection{Systems}

Table~\ref{tab:systems} lists the systems. Every system receives the same bug report and returns at most ten functions, most likely first. Answers given as file paths and names are mapped to the functions at the base commit with Tree-sitter; an answer that matches no function is a miss.

\begin{table}[t]
\centering
\small
\caption{Systems compared. \emph{Free}: usable at no charge by a developer (decided by system, before any result was read). Agents ran on the \NAgent{}-bug subset; the other systems also ran on all \NBugs{} bugs where access allowed.}
\label{tab:systems}
\begin{adjustbox}{max width=\linewidth}\begin{tabular}{llllc}
\toprule
Family & System & Model & Access & Free \\
\midrule
Free localizer & \tool, browser only & learned ranker, static embeddings & runs in the browser & \checkmark \\
 & \toolgpu & + Qwen3-Reranker-0.6B with LoRA & one GPU & \checkmark \\
Retrieval & SweRankEmbed-Small, -Large~\cite{reddy2025swerank} & published checkpoints & one GPU & \checkmark \\
Single call & gpt-6-luna & OpenAI & Codex CLI, shell disabled & \checkmark \\
 & gpt-6-astra & OpenAI & Codex CLI, shell disabled & \\
 & Gemini 3.8 Flash & Google & Gemini API & \checkmark \\
 & Claude Opus 5.5, Sonnet 5.5 & Anthropic & Claude Code, tools disabled & \\
 & Qwen3-Coder-Next (FP8) & open weights & vLLM, 4$\times$L40S & \checkmark \\
 & gpt-oss-120b~\cite{openai2025gptoss} & open weights & vLLM, 2--4$\times$L40S & \checkmark \\
Agent & gpt-6-luna, gpt-6-astra & OpenAI & Codex CLI, read-only sandbox & luna \\
 & Claude Opus 5.5 & Anthropic & Claude Code, read-only tools & \\
 & Qwen3-Coder-Next, gpt-oss-120b & open weights & my agent loop, 40 turns & \checkmark \\
\bottomrule
\end{tabular}\end{adjustbox}
\end{table}

\textbf{\tool.} \tool runs in a web worker. It downloads the repository at the base commit, splits every source file into functions with Tree-sitter compiled to WebAssembly, and indexes them with a BM25 variant that tokenizes identifiers. It embeds every function with a 16.8\,MB static code-embedding model, which needs no neural network at encoding time, and ranks candidates with a two-stage LambdaMART ranker over lexical, structural, mention, and embedding features. The ranker was trained on 7,845 bugs, in all six languages, from projects disjoint from \fresh. No code or report leaves the user's device.

\textbf{\toolgpu.} The same pipeline, with its top 50 candidates re-read by Qwen3-Reranker-0.6B~\cite{zhang2025qwen3embedding} with one LoRA adapter~\cite{hu2022lora} trained on bugs from all six languages, and the reranker's score blended with the browser's. Its configuration was selected on training projects only (Section~\ref{sec:reranker}) and frozen before it ran on any \fresh bug.

\textbf{SweRank.} SweRankEmbed-Small and -Large~\cite{reddy2025swerank} embed the report and every function at the base commit and rank functions by similarity.

\textbf{Single-call models.} The model receives the report and the repository's file list in one request and returns ten functions as JSON. It has no tools and cannot read code. The frontier models run through their vendors' own command-line tools with every tool disabled, or through the vendor's API; the open-weight models are served with vLLM on the same hardware for every bug.

\textbf{Agents.} The model receives the same report, a read-only checkout of the repository, and tools to list, search, and read files, with web access disabled, and returns ten functions. Claude Opus~5.5 runs in Claude Code and the OpenAI models in Codex, as developers use them; Claude Code also delegates some sub-tasks to Claude Haiku~5.5. The open-weight models run in a minimal agent loop of my own with the same three read-only tools and an answer tool, for at most 40 turns.

\textbf{Free to use.} I classify a system as free if a developer can use it at no charge: \tool, SweRank, the open-weight models, Gemini~3.8 Flash, which is available on Google's free tier, and gpt-6-luna, the model in ChatGPT's free plan. This classification was fixed by system before any result was read.

\subsection{Metrics}

The primary metric is function-level hit@10, the share of bugs for which at least one changed function is among the ten returned. I also report hit@1 and hit@5, all@10 (every changed function is among the ten), and mean reciprocal rank (MRR). Cost per bug is measured at list prices on the run date: for model-based systems, the tokens each run used at the provider's API rates, including runs made under a subscription; for GPU systems, GPU time at a public rental rate.

\subsection{Statistical Analysis}

Bugs from the same project are correlated, so every 95\% confidence interval comes from a bootstrap that resamples projects rather than bugs. Paired comparisons use the exact McNemar test~\cite{mcnemar1947}, with Holm's correction~\cite{holm1979} across all comparisons in a table. I report every comparison against both \toolgpu and \tool, each as a separate family. With \NAgent{} bugs and about 20\% of pairs discordant, the subset detects a paired difference of about nine points with power 0.8 at $\alpha=0.05$; smaller differences between top systems are reported as ties, not as evidence of equivalence.

\subsection{Escalation}

\tool estimates its own confidence: a logistic regression on ten features of its ranking, such as the score gap between the first and second candidates, trained on its own out-of-fold training predictions, predicts the probability that a changed function is in its top ten. A bug is escalated to an agent when this probability is below 0.7899. I fixed the threshold before the study as the 40th percentile of the model's predictions on development bugs. Escalated bugs take the agent's answer; the rest keep \tool's. The baseline escalates the same number of bugs at random.

\subsection{Protocol and Deviations}

The protocol and its amendments are in the artifact. Each system gives one answer per bug, in a fixed random order; I measure repeatability by asking gpt-6-astra again on 50 bugs. A run stopped by a usage limit is reported on the prefix it completed, which is itself a random sample because the order is random. The amendments added the full-report condition, fixed gaps in the open-weight agents' harness that I found in run records before reading any answer, replaced models that I could not access, and set the order of rate-limited runs. No amendment was written after reading the results it affects.

%% ---------------------------------------------------------------------------
\section{Results}
\label{sec:results}

\subsection{RQ1: How Accurate Is Each System?}
\label{sec:rq1}

\begin{table*}[t]
\centering
\caption{Full bug reports, the 198-bug agent subset: every free system and \tool{}. hit@k: a function the fix changed is in the top $k$; all@10: every changed function is in the top 10. 95\% intervals from a bootstrap over projects. $\Delta$: against \toolgpu{}, with wins/losses and the exact paired test, Holm-adjusted across the table.}
\label{tab:full198}
\begin{tabular}{lcrrrrrrrr}
\toprule
System & Free & $n$ & hit@1 & hit@10 & 95\% CI & all@10 & MRR & Cost/bug & $\Delta$ hit@10 (W/L, $p_{\mathrm{Holm}}$) \\
\midrule
gpt-6-luna, agent & \checkmark & 198 & 65.2 & 77.3 & 70.9--83.3 & 36.4 & 0.702 & \$0.003 & +0.5 (26/25, 1.000) \\
\tool{}, browser only & \checkmark & 198 & 46.0 & 77.3 & 71.0--83.1 & 42.9 & 0.566 & \$0 & +0.5 (7/6, 1.000) \\
\textbf{\toolgpu{}} & \checkmark & 198 & 47.5 & 76.8 & 69.8--83.0 & 45.5 & 0.566 & \$0.001 & reference \\
SweRankEmbed-Large & \checkmark & 198 & 40.4 & 72.7 & 66.1--79.6 & 38.4 & 0.518 & -- & -4.0 (14/22, 1.000) \\
gpt-6-astra, one call &  & 95 & 55.8 & 67.4 & 54.8--78.8 & 28.4 & 0.601 & \$0.62 & -10.5 (8/18, 1.000) \\
SweRankEmbed-Small & \checkmark & 198 & 38.4 & 62.1 & 54.5--69.7 & 33.8 & 0.459 & -- & -14.6 (11/40, 0.004) \\
Qwen3-Coder-Next, agent & \checkmark & 198 & 53.5 & 61.6 & 53.5--69.5 & 34.8 & 0.566 & \$0.027 & -15.2 (15/45, 0.004) \\
Gemini 3.8 Flash, one call & \checkmark & 198 & 42.4 & 55.6 & 46.3--64.9 & 28.3 & 0.468 & \$0.100 & -21.2 (11/53, $<$0.001) \\
gpt-oss-120b, agent & \checkmark & 198 & 48.5 & 51.0 & 42.3--59.5 & 23.7 & 0.497 & \$0.049 & -25.8 (10/61, $<$0.001) \\
gpt-6-luna, one call & \checkmark & 198 & 37.4 & 42.9 & 33.7--52.4 & 20.2 & 0.397 & \$0.005 & -33.8 (4/71, $<$0.001) \\
Qwen3-Coder-Next, one call & \checkmark & 198 & 31.8 & 40.9 & 31.7--50.0 & 20.7 & 0.351 & \$0.006 & -35.9 (5/76, $<$0.001) \\
gpt-oss-120b, one call & \checkmark & 198 & 25.8 & 31.3 & 22.8--40.0 & 14.1 & 0.277 & \$0.006 & -45.5 (3/93, $<$0.001) \\
\bottomrule
\end{tabular}
\end{table*}


\textbf{The agent subset.} Table~\ref{tab:full198} compares every system on the \NAgent{}-bug subset with full reports. \tool reaches \res{\FF}{website}{hit@10}\% hit@10 in the browser alone and \res{\FF}{website-gpu}{hit@10}\% with the GPU reranker, a difference within noise on this subset. The best free agent, gpt-6-luna in Codex, reaches \res{\FF}{codex-agent}{hit@10}\%. It wins on \res{\FF}{codex-agent}{wlhit@10} bugs against \toolgpu (wins/losses), so the two are tied. Every other free system is less accurate. SweRankEmbed-Large trails \toolgpu by \res{\FF}{swerank-embed-large}{dhit@10} points, which is not significant on this subset. The open-weight agents trail by \res{\FF}{open-agent-qwen3-coder-next-fp8}{dhit@10} (Qwen3-Coder-Next) and \res{\FF}{open-agent-gpt-oss-120b}{dhit@10} points (gpt-oss-120b), and every single-call model by more than 20 points, all with $p_\mathrm{Holm}<0.01$.

\begin{figure}[t]
\centering
\begin{tikzpicture}
\begin{axis}[width=0.7\linewidth, height=6cm, xtick={1,5,10}, xmin=0.5, xmax=10.5, ymin=10, ymax=90,
  xlabel={$k$}, ylabel={Function hit@$k$ (\%)}, grid=major, grid style={gray!20}, tick label style={font=\footnotesize},
  label style={font=\small}, legend style={font=\footnotesize, at={(1.03,0.5)}, anchor=west}]
\addplot[mark=*, color=black] coordinates {(1,47.47) (5,69.70) (10,76.77)};
\addlegendentry{\toolgpu{}}
\addplot[mark=o, color=black] coordinates {(1,45.96) (5,69.19) (10,77.27)};
\addlegendentry{\tool{}, browser only}
\addplot[mark=triangle*, color=black] coordinates {(1,65.15) (5,77.27) (10,77.27)};
\addlegendentry{gpt-6-luna, agent}
\addplot[mark=square*, color=black] coordinates {(1,40.40) (5,65.15) (10,72.73)};
\addlegendentry{SweRankEmbed-Large}
\addplot[mark=diamond*, color=black] coordinates {(1,53.54) (5,61.62) (10,61.62)};
\addlegendentry{Qwen3-Coder-Next, agent}
\addplot[mark=triangle, color=black] coordinates {(1,42.42) (5,53.03) (10,55.56)};
\addlegendentry{Gemini 3.8 Flash, one call}
\addplot[mark=square, color=black] coordinates {(1,37.37) (5,42.93) (10,42.93)};
\addlegendentry{gpt-6-luna, one call}
\end{axis}
\end{tikzpicture}
\caption{Function hit@$k$ for $k=1,5,10$, full reports, 198-bug subset.}
\label{fig:hitk}
\end{figure}


\textbf{Ranking quality.} The agent ranks better than \tool. Its hit@1 is \res{\FF}{codex-agent}{hit@1}\%, against \res{\FF}{website-gpu}{hit@1}\% for \toolgpu ($p_\mathrm{Holm}\res{\FF}{codex-agent}{phit@1}$): when the agent finds a faulty function, it usually puts it first (Figure~\ref{fig:hitk}). \tool more often has every changed function in its top ten (all@10 of \res{\FF}{website-gpu}{all@10}\% against \res{\FF}{codex-agent}{all@10}\%), which matters for the many fixes that change more than one function.

\begin{table*}[t]
\centering
\caption{Full bug reports, all 580 bugs: the systems that ran on the whole set. Columns as in Table~\ref{tab:full198}.}
\label{tab:full580}
\begin{adjustbox}{max width=\linewidth}\begin{tabular}{lcrrrrrrrr}
\toprule
System & Free & $n$ & hit@1 & hit@10 & 95\% CI & all@10 & MRR & Cost/bug & $\Delta$ hit@10 (W/L, $p_{\mathrm{Holm}}$) \\
\midrule
\textbf{\toolgpu{}} & \checkmark & 580 & 46.7 & 77.9 & 73.4--82.4 & 45.5 & 0.575 & \$0.001 & reference \\
\tool{}, browser only & \checkmark & 580 & 46.0 & 75.3 & 70.1--80.1 & 41.4 & 0.560 & \$0 & -2.6 (12/27, 0.166) \\
SweRankEmbed-Large & \checkmark & 580 & 43.3 & 72.9 & 66.5--78.5 & 37.9 & 0.538 & -- & -5.0 (35/64, 0.046) \\
SweRankEmbed-Small & \checkmark & 580 & 37.4 & 65.9 & 59.4--71.9 & 35.0 & 0.470 & -- & -12.1 (30/100, $<$0.001) \\
gpt-6-astra, one call &  & 372 & 44.4 & 54.6 & 46.7--62.0 & 25.5 & 0.485 & \$0.46 & -24.2 (21/111, $<$0.001) \\
Gemini 3.8 Flash, one call & \checkmark & 398 & 41.5 & 53.5 & 45.9--60.5 & 27.6 & 0.452 & \$0.11 & -25.6 (19/121, $<$0.001) \\
gpt-6-luna, one call & \checkmark & 580 & 34.0 & 39.8 & 33.3--46.3 & 19.1 & 0.366 & \$0.005 & -38.1 (8/229, $<$0.001) \\
Qwen3-Coder-Next, one call & \checkmark & 580 & 29.0 & 36.7 & 30.6--42.8 & 18.3 & 0.318 & \$0.006 & -41.2 (10/249, $<$0.001) \\
gpt-oss-120b, one call & \checkmark & 580 & 21.7 & 27.2 & 21.1--33.5 & 13.6 & 0.239 & \$0.008 & -50.7 (4/298, $<$0.001) \\
\bottomrule
\end{tabular}\end{adjustbox}
\end{table*}


\textbf{All \NBugs{} bugs.} Table~\ref{tab:full580} covers the systems that ran on the whole benchmark. \toolgpu reaches \res{\FFA}{website-gpu}{hit@10}\% (95\% CI \res{\FFA}{website-gpu}{hit@10-ci}). SweRankEmbed-Large trails it by \res{\FFA}{swerank-embed-large}{dhit@10} points ($p_\mathrm{Holm}\res{\FFA}{swerank-embed-large}{phit@10}$) and gpt-6-luna asked once by \res{\FFA}{codex-single}{dhit@10}. The GPU reranker adds \res{\FFA}{website-gpu}{browser-dhit@10} points over the browser alone, which is not significant after correction for the other comparisons in its family ($p_\mathrm{Holm}\res{\FFA}{website-gpu}{browser-phit@10}$).

\textbf{Paid frontier agents.} With full reports, Claude Opus~5.5 as an agent reaches \res{\FF}{claude-agent}{hit@10}\% hit@10 and gpt-6-astra \res{\FF}{codex-agent-astra}{hit@10}\%; gpt-6-astra asked once reaches \res{\FF}{codex-single-astra}{hit@10}\% on the \res{\FF}{codex-single-astra}{n} bugs completed so far. On shortened reports, where both agents completed every bug, they were the two most accurate systems in the study (Table~\ref{tab:short198}).

\begin{table}[t]
\centering
\small
\caption{Bugs for which a system returned no function, full reports, 198-bug subset. Each counts as a miss. \emph{Turns}: the agent used all 40 turns and a final request without answering; \emph{Context}: the prompt or conversation outgrew the model's context; \emph{Empty}: the reply had no function that could be read.}
\label{tab:noanswer}
\begin{adjustbox}{max width=\linewidth}\begin{tabular}{lrrrrr}
\toprule
System & Bugs & Turns & Context & Empty & Total \\
\midrule
\toolgpu{} & 198 & 0 & 0 & 0 & 0 \\
\tool{}, browser only & 198 & 0 & 0 & 0 & 0 \\
gpt-6-luna, agent & 198 & 0 & 0 & 2 & 2 \\
gpt-6-luna, one call & 198 & 0 & 0 & 10 & 10 \\
Gemini 3.8 Flash, one call & 198 & 0 & 0 & 0 & 0 \\
Qwen3-Coder-Next, agent & 198 & 12 & 0 & 17 & 29 \\
Qwen3-Coder-Next, one call & 198 & 0 & 5 & 7 & 12 \\
gpt-oss-120b, agent & 198 & 36 & 7 & 5 & 48 \\
gpt-oss-120b, one call & 198 & 0 & 32 & 1 & 33 \\
gpt-6-astra, agent & 198 & 0 & 0 & 2 & 2 \\
gpt-6-astra, one call & 198 & 0 & 0 & 9 & 9 \\
SweRankEmbed-Large & 198 & 0 & 0 & 0 & 0 \\
SweRankEmbed-Small & 198 & 0 & 0 & 0 & 0 \\
\bottomrule
\end{tabular}\end{adjustbox}
\end{table}


\textbf{Unanswered bugs.} Some systems returned no usable answer for some bugs, and the protocol counts these as misses (Table~\ref{tab:noanswer}). The open-weight agents account for most of them: gpt-oss-120b used all 40 turns without answering on many bugs and outgrew its context on others, and some full reports do not fit in its 131,072-token context even as a single call. These failures are part of what these models are as localizers today, but they also mean that a different harness could score them higher.

\begin{finding}{Answer to RQ1}
On unseen bugs, the free browser localizer is tied with the best free coding agent at hit@10 and is significantly more accurate than SweRank on the full benchmark, both open-weight models, and every single-call model. The free agent ranks the faulty function first far more often. Paid frontier agents are the most accurate systems overall.
\end{finding}

\subsection{RQ2: What Does Accuracy Cost?}
\label{sec:rq2}

\begin{figure}[t]
\centering
\begin{tikzpicture}
\begin{axis}[width=0.92\linewidth, height=6.2cm, xmode=log, xmin=5e-5, xmax=2,
  xlabel={Cost per bug (US\$, log scale; free systems at the left edge)}, ylabel={Function hit@10 (\%)},
  ymin=20, ymax=100, grid=major, grid style={gray!20}, tick label style={font=\footnotesize},
  label style={font=\small}, legend style={font=\footnotesize, at={(0.98,0.03)}, anchor=south east}]
\addplot[only marks, mark=*, color=black, mark size=2.4pt] coordinates {(0.001178,76.77) (0.000100,77.27) (0.003188,77.27) (0.005300,42.93) (0.099581,55.56) (0.026891,61.62) (0.006490,40.91) (0.049141,51.01) (0.006322,31.31)};
\addlegendentry{free to use}
\addplot[only marks, mark=triangle*, color=gray, mark size=2.4pt] coordinates {(0.535998,67.37)};
\addlegendentry{paid}
\node[font=\scriptsize, anchor=west] at (axis cs:0.001472,76.77) {\toolgpu{}};
\node[font=\scriptsize, anchor=west] at (axis cs:0.000125,77.27) {\tool{}, browser only};
\node[font=\scriptsize, anchor=west] at (axis cs:0.003985,77.27) {gpt-6-luna, agent};
\node[font=\scriptsize, anchor=west] at (axis cs:0.006625,42.93) {gpt-6-luna, one call};
\node[font=\scriptsize, anchor=west] at (axis cs:0.124476,55.56) {Gemini 3.8 Flash, one call};
\node[font=\scriptsize, anchor=west] at (axis cs:0.033613,61.62) {Qwen3-Coder-Next, agent};
\node[font=\scriptsize, anchor=west] at (axis cs:0.008113,40.91) {Qwen3-Coder-Next, one call};
\node[font=\scriptsize, anchor=west] at (axis cs:0.061427,51.01) {gpt-oss-120b, agent};
\node[font=\scriptsize, anchor=west] at (axis cs:0.007902,31.31) {gpt-oss-120b, one call};
\node[font=\scriptsize, anchor=west] at (axis cs:0.669998,67.37) {gpt-6-astra, one call};
\end{axis}
\end{tikzpicture}
\caption{Accuracy against cost per bug, full reports, 198-bug subset. Costs are list prices for the tokens each run used, or GPU rental time; systems that cost nothing to run sit at the left edge.}
\label{fig:cost}
\end{figure}


Figure~\ref{fig:cost} plots accuracy against cost per bug. \tool in the browser costs nothing to run, and the GPU reranker costs about \CostAFgpu{} per bug in GPU time. gpt-6-luna as an agent costs \res{\FF}{codex-agent}{cost} per bug at API prices, gpt-6-astra as an agent \CostAstraAgent{}, and Claude Opus~5.5 as an agent \CostClaudeAgent{}. Among the systems that ran on all \NBugs{} bugs, the two \tool variants are the only ones that no other system outperforms on both accuracy and cost.

Speed and privacy differ as well. \tool's median time per bug is \res{\FF}{website}{sec}\,s in the browser, and the agents' is \res{\FF}{codex-agent}{sec}\,s for gpt-6-luna. More importantly for many organizations, every model-based system sends the repository and the report to a third party; \tool sends nothing.

\begin{finding}{Answer to RQ2}
The free localizer is on the cost--accuracy frontier: it costs nothing in the browser and about a tenth of a cent with the GPU reranker, and it is the only option that keeps code on the developer's machine. Paid frontier agents buy higher accuracy at about two hundred times the cost of the GPU reranker.
\end{finding}

\subsection{RQ3: Can Escalation Get the Best of Both?}
\label{sec:rq3}

\begin{table}[t]
\centering
\small
\caption{Confidence-gated escalation (RQ3): the bugs \tool{} is least confident about, at the threshold fixed before the study (0.7899), go to an agent; the rest keep \tool{}'s answer. Random escalates the same share at random.}
\label{tab:escalation}
\resizebox{\columnwidth}{!}{\begin{tabular}{lrrrrr}
\toprule
Agent & Escalated & Cascade & Random & Agent alone & Cascade cost/bug \\
\midrule
gpt-6-luna, agent & 56.1\% & 81.8 & 77.0 & 77.3 & \$0.003 \\
Qwen3-Coder-Next, agent & 56.1\% & 67.2 & 68.3 & 61.6 & \$0.016 \\
gpt-oss-120b, agent & 56.1\% & 63.6 & 62.3 & 51.0 & \$0.029 \\
\bottomrule
\end{tabular}}
\end{table}

\begin{figure}[t]
\centering
\begin{tikzpicture}
\begin{axis}[width=0.49\linewidth, height=5.4cm, title={\small Full reports, gpt-6-luna agent}, xmin=0, xmax=100, ymin=55, ymax=100,
  xlabel={Bugs escalated (\%)}, ylabel={Function hit@10 (\%)}, grid=major, grid style={gray!20},
  tick label style={font=\footnotesize}, label style={font=\small}, legend style={font=\footnotesize, at={(0.98,0.03)}, anchor=south east}]
\addplot[thick, color=black, mark=none] coordinates {(0.00,76.77) (0.00,76.77) (0.51,76.77) (1.01,77.27) (1.52,77.27) (2.02,77.27) (2.53,77.78) (3.03,77.27) (3.54,77.78) (4.04,77.27) (4.55,77.27) (5.05,77.27) (5.56,77.27) (6.06,77.78) (6.57,78.28) (7.07,78.28) (7.58,78.28) (8.08,78.28) (8.59,77.78) (9.09,77.78) (9.60,77.78) (10.10,77.27) (10.61,77.27) (11.11,77.27) (11.62,76.77) (12.12,77.27) (12.63,77.27) (13.13,77.27) (13.64,77.27) (14.14,77.78) (14.65,77.78) (15.15,77.78) (15.66,77.78) (16.16,78.28) (16.67,78.79) (17.17,79.29) (17.68,79.29) (18.18,78.79) (18.69,79.29) (19.19,79.29) (19.70,79.29) (20.20,79.80) (20.71,79.80) (21.21,79.80) (21.72,79.80) (22.22,80.30) (22.73,79.80) (23.23,79.29) (23.74,79.80) (24.24,79.29) (24.75,79.80) (25.25,79.80) (25.76,79.80) (26.26,79.29) (26.77,79.29) (27.27,79.80) (27.78,79.80) (28.28,80.30) (28.79,80.30) (29.29,80.30) (29.80,80.30) (30.30,80.30) (30.81,80.30) (31.82,80.81) (32.32,80.81) (32.83,80.30) (33.33,80.30) (33.84,80.30) (34.34,80.30) (34.85,80.81) (35.35,80.81) (35.86,80.81) (36.36,80.81) (36.87,80.81) (37.37,81.31) (37.88,81.31) (38.38,81.31) (38.89,81.31) (39.39,81.31) (39.90,81.31) (40.40,81.31) (40.91,81.31) (41.41,81.31) (41.92,81.31) (42.42,81.31) (42.93,81.31) (43.43,81.31) (43.94,81.31) (44.44,81.82) (44.95,81.82) (45.45,82.32) (45.96,82.32) (46.46,82.32) (46.97,82.32) (47.47,82.32) (47.98,82.32) (48.48,82.32) (48.99,82.32) (49.49,82.32) (50.00,82.32) (50.51,81.82) (51.01,81.82) (51.52,81.82) (52.02,82.32) (52.53,81.82) (53.03,81.82) (53.54,81.82) (54.04,81.31) (54.55,81.82) (55.05,81.82) (55.56,81.82) (56.06,81.82) (56.57,81.82) (57.07,81.82) (57.58,81.82) (58.08,81.82) (58.59,81.82) (59.09,81.82) (59.60,81.82) (60.10,81.31) (60.61,81.31) (61.11,80.81) (61.62,80.81) (62.12,80.81) (62.63,80.81) (63.13,81.31) (63.64,81.31) (64.14,81.31) (64.65,81.31) (65.15,80.81) (65.66,80.81) (66.16,80.81) (66.67,80.81) (67.17,80.30) (67.68,80.30) (68.18,80.30) (68.69,79.80) (69.19,79.80) (69.70,79.80) (70.20,79.80) (70.71,79.80) (71.21,79.80) (71.72,79.80) (72.22,79.29) (72.73,79.29) (73.23,79.29) (73.74,79.29) (74.24,79.29) (74.75,79.29) (75.25,79.80) (75.76,79.80) (76.26,79.80) (77.27,79.29) (77.78,79.29) (78.28,79.29) (78.79,79.29) (79.29,79.29) (79.80,79.29) (80.30,79.29) (80.81,79.29) (81.31,79.29) (81.82,78.79) (82.32,78.79) (82.83,78.28) (83.33,78.28) (83.84,78.28) (84.85,77.78) (85.35,77.78) (85.86,77.78) (86.36,77.78) (86.87,77.78) (87.37,77.78) (87.88,77.78) (88.38,77.78) (88.89,77.78) (89.90,77.78) (90.40,77.78) (90.91,77.78) (91.41,77.78) (91.92,77.78) (92.42,77.27) (92.93,77.27) (93.43,77.27) (93.94,77.27) (94.44,77.27) (94.95,77.27) (95.45,77.27) (95.96,77.27) (96.46,77.27) (96.97,77.27) (97.47,77.27) (97.98,77.27) (98.48,77.27) (98.99,77.27) (99.49,77.27) (100.00,77.27)};
\addlegendentry{confidence gate}
\addplot[thick, dashed, color=gray, mark=none] coordinates {(0.00,76.77) (0.00,76.77) (0.51,76.77) (1.01,76.77) (1.52,76.78) (2.02,76.78) (2.53,76.78) (3.03,76.78) (3.54,76.79) (4.04,76.79) (4.55,76.79) (5.05,76.79) (5.56,76.80) (6.06,76.80) (6.57,76.80) (7.07,76.80) (7.58,76.81) (8.08,76.81) (8.59,76.81) (9.09,76.81) (9.60,76.82) (10.10,76.82) (10.61,76.82) (11.11,76.82) (11.62,76.83) (12.12,76.83) (12.63,76.83) (13.13,76.83) (13.64,76.84) (14.14,76.84) (14.65,76.84) (15.15,76.84) (15.66,76.85) (16.16,76.85) (16.67,76.85) (17.17,76.85) (17.68,76.86) (18.18,76.86) (18.69,76.86) (19.19,76.86) (19.70,76.87) (20.20,76.87) (20.71,76.87) (21.21,76.87) (21.72,76.88) (22.22,76.88) (22.73,76.88) (23.23,76.89) (23.74,76.89) (24.24,76.89) (24.75,76.89) (25.25,76.90) (25.76,76.90) (26.26,76.90) (26.77,76.90) (27.27,76.91) (27.78,76.91) (28.28,76.91) (28.79,76.91) (29.29,76.92) (29.80,76.92) (30.30,76.92) (30.81,76.92) (31.82,76.93) (32.32,76.93) (32.83,76.93) (33.33,76.94) (33.84,76.94) (34.34,76.94) (34.85,76.94) (35.35,76.95) (35.86,76.95) (36.36,76.95) (36.87,76.95) (37.37,76.96) (37.88,76.96) (38.38,76.96) (38.89,76.96) (39.39,76.97) (39.90,76.97) (40.40,76.97) (40.91,76.97) (41.41,76.98) (41.92,76.98) (42.42,76.98) (42.93,76.98) (43.43,76.99) (43.94,76.99) (44.44,76.99) (44.95,76.99) (45.45,77.00) (45.96,77.00) (46.46,77.00) (46.97,77.00) (47.47,77.01) (47.98,77.01) (48.48,77.01) (48.99,77.02) (49.49,77.02) (50.00,77.02) (50.51,77.02) (51.01,77.03) (51.52,77.03) (52.02,77.03) (52.53,77.03) (53.03,77.04) (53.54,77.04) (54.04,77.04) (54.55,77.04) (55.05,77.05) (55.56,77.05) (56.06,77.05) (56.57,77.05) (57.07,77.06) (57.58,77.06) (58.08,77.06) (58.59,77.06) (59.09,77.07) (59.60,77.07) (60.10,77.07) (60.61,77.07) (61.11,77.08) (61.62,77.08) (62.12,77.08) (62.63,77.08) (63.13,77.09) (63.64,77.09) (64.14,77.09) (64.65,77.09) (65.15,77.10) (65.66,77.10) (66.16,77.10) (66.67,77.10) (67.17,77.11) (67.68,77.11) (68.18,77.11) (68.69,77.11) (69.19,77.12) (69.70,77.12) (70.20,77.12) (70.71,77.12) (71.21,77.13) (71.72,77.13) (72.22,77.13) (72.73,77.13) (73.23,77.14) (73.74,77.14) (74.24,77.14) (74.75,77.15) (75.25,77.15) (75.76,77.15) (76.26,77.15) (77.27,77.16) (77.78,77.16) (78.28,77.16) (78.79,77.17) (79.29,77.17) (79.80,77.17) (80.30,77.17) (80.81,77.18) (81.31,77.18) (81.82,77.18) (82.32,77.18) (82.83,77.19) (83.33,77.19) (83.84,77.19) (84.85,77.20) (85.35,77.20) (85.86,77.20) (86.36,77.20) (86.87,77.21) (87.37,77.21) (87.88,77.21) (88.38,77.21) (88.89,77.22) (89.90,77.22) (90.40,77.22) (90.91,77.23) (91.41,77.23) (91.92,77.23) (92.42,77.23) (92.93,77.24) (93.43,77.24) (93.94,77.24) (94.44,77.24) (94.95,77.25) (95.45,77.25) (95.96,77.25) (96.46,77.25) (96.97,77.26) (97.47,77.26) (97.98,77.26) (98.48,77.27) (98.99,77.27) (99.49,77.27) (100.00,77.27)};
\addlegendentry{random}
\addplot[only marks, mark=*, mark size=2.6pt, color=black, forget plot] coordinates {(56.06,81.82)};
\node[font=\scriptsize, anchor=south east] at (axis cs:56.06,82.82) {fixed threshold};
\end{axis}
\end{tikzpicture}\hfill\begin{tikzpicture}
\begin{axis}[width=0.49\linewidth, height=5.4cm, title={\small Shortened reports, Claude Opus 5.5 agent}, xmin=0, xmax=100, ymin=55, ymax=100,
  xlabel={Bugs escalated (\%)}, ylabel={Function hit@10 (\%)}, grid=major, grid style={gray!20},
  tick label style={font=\footnotesize}, label style={font=\small}, legend style={font=\footnotesize, at={(0.98,0.03)}, anchor=south east}]
\addplot[thick, color=black, mark=none] coordinates {(0.00,74.24) (0.00,74.24) (0.51,74.75) (1.01,74.75) (1.52,74.75) (2.02,74.75) (2.53,75.25) (3.03,74.75) (3.54,75.25) (4.04,75.76) (4.55,75.76) (5.05,75.76) (5.56,75.76) (6.06,75.25) (6.57,75.25) (7.07,75.25) (7.58,75.76) (8.59,76.77) (9.09,77.27) (9.60,77.27) (10.10,77.27) (10.61,77.78) (11.11,77.78) (11.62,78.28) (12.12,78.28) (12.63,78.28) (13.13,78.28) (13.64,78.28) (14.14,78.28) (14.65,78.28) (15.15,78.79) (15.66,79.29) (16.16,79.80) (16.67,79.80) (17.17,79.80) (17.68,79.80) (18.18,79.80) (18.69,79.80) (19.19,79.80) (19.70,79.80) (20.20,79.80) (20.71,80.30) (21.21,80.81) (21.72,80.81) (22.22,81.31) (22.73,81.31) (23.23,81.31) (23.74,81.31) (24.24,80.81) (24.75,80.81) (25.25,81.31) (25.76,81.82) (26.26,81.82) (26.77,82.32) (27.27,82.32) (27.78,82.32) (28.28,82.83) (28.79,82.83) (29.29,82.83) (29.80,82.83) (30.30,82.83) (30.81,82.83) (31.31,83.33) (31.82,83.84) (32.32,83.84) (32.83,83.84) (33.33,84.34) (33.84,84.34) (34.34,84.85) (34.85,84.85) (35.35,85.35) (35.86,85.35) (36.36,85.86) (36.87,86.36) (37.37,86.87) (37.88,87.37) (38.38,87.37) (38.89,87.37) (39.39,87.88) (39.90,88.38) (40.40,88.38) (40.91,88.38) (41.41,88.38) (41.92,88.89) (42.42,88.89) (42.93,88.89) (43.43,88.89) (43.94,89.39) (44.44,89.90) (44.95,89.90) (45.45,89.90) (45.96,89.90) (46.46,89.90) (46.97,89.90) (47.47,89.90) (47.98,89.90) (48.48,89.90) (48.99,90.40) (49.49,90.91) (50.51,90.91) (51.01,90.91) (51.52,90.91) (52.02,90.91) (52.53,90.91) (53.03,90.91) (53.54,90.91) (54.04,90.91) (54.55,90.91) (55.05,90.91) (55.56,90.91) (56.06,90.91) (56.57,91.41) (57.07,91.41) (57.58,91.41) (58.08,90.91) (58.59,90.91) (59.09,90.91) (59.60,90.91) (60.10,90.91) (60.61,90.91) (61.11,90.91) (61.62,91.41) (62.12,91.41) (62.63,91.41) (63.13,91.41) (63.64,91.92) (64.14,91.92) (64.65,91.92) (65.15,91.92) (65.66,92.42) (66.16,92.42) (66.67,91.92) (67.17,91.92) (67.68,91.92) (68.18,91.92) (68.69,91.92) (69.19,91.92) (69.70,91.92) (70.20,91.92) (70.71,91.92) (71.21,91.92) (71.72,91.92) (72.22,91.92) (72.73,91.92) (73.74,91.92) (74.24,91.92) (74.75,91.92) (75.25,91.41) (75.76,91.41) (76.26,91.41) (76.77,91.41) (77.27,91.41) (77.78,91.41) (78.28,91.41) (78.79,91.92) (79.29,91.92) (80.30,91.92) (80.81,91.92) (81.31,91.92) (81.82,91.92) (82.32,91.92) (82.83,91.92) (83.33,91.92) (83.84,91.92) (84.34,91.92) (84.85,91.92) (85.35,91.92) (85.86,91.92) (86.36,91.92) (86.87,91.92) (87.37,91.41) (87.88,91.41) (88.38,91.41) (88.89,91.41) (89.39,91.41) (89.90,91.41) (90.40,91.41) (90.91,91.41) (91.41,91.41) (91.92,90.91) (92.42,90.91) (92.93,90.91) (93.43,90.91) (93.94,90.91) (94.44,90.91) (94.95,90.91) (95.45,90.91) (95.96,90.91) (96.46,90.91) (96.97,90.91) (97.47,90.91) (97.98,90.91) (98.48,90.91) (98.99,90.91) (99.49,90.91) (100.00,90.91)};
\addlegendentry{confidence gate}
\addplot[thick, dashed, color=gray, mark=none] coordinates {(0.00,74.24) (0.00,74.24) (0.51,74.33) (1.01,74.41) (1.52,74.49) (2.02,74.58) (2.53,74.66) (3.03,74.75) (3.54,74.83) (4.04,74.92) (4.55,75.00) (5.05,75.08) (5.56,75.17) (6.06,75.25) (6.57,75.34) (7.07,75.42) (7.58,75.51) (8.59,75.67) (9.09,75.76) (9.60,75.84) (10.10,75.93) (10.61,76.01) (11.11,76.09) (11.62,76.18) (12.12,76.26) (12.63,76.35) (13.13,76.43) (13.64,76.52) (14.14,76.60) (14.65,76.68) (15.15,76.77) (15.66,76.85) (16.16,76.94) (16.67,77.02) (17.17,77.10) (17.68,77.19) (18.18,77.27) (18.69,77.36) (19.19,77.44) (19.70,77.53) (20.20,77.61) (20.71,77.69) (21.21,77.78) (21.72,77.86) (22.22,77.95) (22.73,78.03) (23.23,78.11) (23.74,78.20) (24.24,78.28) (24.75,78.37) (25.25,78.45) (25.76,78.54) (26.26,78.62) (26.77,78.70) (27.27,78.79) (27.78,78.87) (28.28,78.96) (28.79,79.04) (29.29,79.12) (29.80,79.21) (30.30,79.29) (30.81,79.38) (31.31,79.46) (31.82,79.55) (32.32,79.63) (32.83,79.71) (33.33,79.80) (33.84,79.88) (34.34,79.97) (34.85,80.05) (35.35,80.13) (35.86,80.22) (36.36,80.30) (36.87,80.39) (37.37,80.47) (37.88,80.56) (38.38,80.64) (38.89,80.72) (39.39,80.81) (39.90,80.89) (40.40,80.98) (40.91,81.06) (41.41,81.14) (41.92,81.23) (42.42,81.31) (42.93,81.40) (43.43,81.48) (43.94,81.57) (44.44,81.65) (44.95,81.73) (45.45,81.82) (45.96,81.90) (46.46,81.99) (46.97,82.07) (47.47,82.15) (47.98,82.24) (48.48,82.32) (48.99,82.41) (49.49,82.49) (50.51,82.66) (51.01,82.74) (51.52,82.83) (52.02,82.91) (52.53,83.00) (53.03,83.08) (53.54,83.16) (54.04,83.25) (54.55,83.33) (55.05,83.42) (55.56,83.50) (56.06,83.59) (56.57,83.67) (57.07,83.75) (57.58,83.84) (58.08,83.92) (58.59,84.01) (59.09,84.09) (59.60,84.18) (60.10,84.26) (60.61,84.34) (61.11,84.43) (61.62,84.51) (62.12,84.60) (62.63,84.68) (63.13,84.76) (63.64,84.85) (64.14,84.93) (64.65,85.02) (65.15,85.10) (65.66,85.19) (66.16,85.27) (66.67,85.35) (67.17,85.44) (67.68,85.52) (68.18,85.61) (68.69,85.69) (69.19,85.77) (69.70,85.86) (70.20,85.94) (70.71,86.03) (71.21,86.11) (71.72,86.20) (72.22,86.28) (72.73,86.36) (73.74,86.53) (74.24,86.62) (74.75,86.70) (75.25,86.78) (75.76,86.87) (76.26,86.95) (76.77,87.04) (77.27,87.12) (77.78,87.21) (78.28,87.29) (78.79,87.37) (79.29,87.46) (80.30,87.63) (80.81,87.71) (81.31,87.79) (81.82,87.88) (82.32,87.96) (82.83,88.05) (83.33,88.13) (83.84,88.22) (84.34,88.30) (84.85,88.38) (85.35,88.47) (85.86,88.55) (86.36,88.64) (86.87,88.72) (87.37,88.80) (87.88,88.89) (88.38,88.97) (88.89,89.06) (89.39,89.14) (89.90,89.23) (90.40,89.31) (90.91,89.39) (91.41,89.48) (91.92,89.56) (92.42,89.65) (92.93,89.73) (93.43,89.81) (93.94,89.90) (94.44,89.98) (94.95,90.07) (95.45,90.15) (95.96,90.24) (96.46,90.32) (96.97,90.40) (97.47,90.49) (97.98,90.57) (98.48,90.66) (98.99,90.74) (99.49,90.82) (100.00,90.91)};
\addlegendentry{random}
\addplot[only marks, mark=*, mark size=2.6pt, color=black, forget plot] coordinates {(62.63,91.41)};
\node[font=\scriptsize, anchor=south east] at (axis cs:62.63,92.41) {fixed threshold};
\end{axis}
\end{tikzpicture}
\caption{Escalating the bugs 	ool{} is least confident about to an agent (solid) against escalating the same number of bugs at random (dashed). The dot is the threshold fixed before the study.}
\label{fig:escalation}
\end{figure}


Table~\ref{tab:escalation} and Figure~\ref{fig:escalation} send the bugs below the fixed confidence threshold to an agent and keep \tool's answer for the rest. With full reports, the threshold escalates \res{\FF}{codex-agent}{esc-escalated}\% of bugs. Escalating them to gpt-6-luna raises hit@10 to \res{\FF}{codex-agent}{esc-hit@10}\%, above both systems alone and above the \res{\FF}{codex-agent}{esc-random_hit@10}\% obtained by escalating the same number of bugs at random. The confidence gate therefore selects the bugs on which the agent's help is worth most.

With a paid agent the gain is larger. On shortened reports, escalating \CascadeClaudeEscalated\% of bugs to Claude Opus~5.5 reaches \CascadeClaude\% hit@10, against \CascadeClaudeRandom\% for random escalation and \ClaudeAgentShort\% for the agent on every bug (Table~\ref{tab:escalation-short}), at roughly two thirds of the agent's cost.

\begin{table}[t]
\centering
\small
\caption{Confidence-gated escalation (RQ3): the bugs Allfather is least confident about, at the threshold fixed before the study (0.7899), go to an agent; the rest keep Allfather's answer. Random escalates the same share at random.}
\label{tab:escalation-short}
\begin{tabular}{lrrrrr}
\toprule
Agent & Escalated & Cascade & Random & Agent alone & Cascade cost/bug \\
\midrule
Claude Opus 5.5, agent & 62.6\% & 91.4 & 84.7 & 90.9 & \$0.13 \\
gpt-6-astra, agent & 62.6\% & 86.9 & 81.8 & 86.4 & \$0.18 \\
gpt-6-luna, agent & 62.6\% & 78.3 & 75.2 & 75.8 & \$0.004 \\
Qwen3-Coder-Next, agent & 62.6\% & 65.7 & 65.7 & 60.6 & \$0.019 \\
\bottomrule
\end{tabular}
\end{table}


\begin{finding}{Answer to RQ3}
Yes. A confidence threshold fixed before the study escalates the bugs on which an agent helps most: with either a free or a paid agent, the cascade beats random escalation of the same size, and with a paid agent it matches the agent's accuracy on every bug at lower cost.
\end{finding}

\subsection{RQ4: Where Do the Systems Differ?}
\label{sec:rq4}

\begin{table}[t]
\centering
\small
\caption{Function hit@10 by language, full reports, all 580 bugs. Bold: best in the column.}
\label{tab:languages}
\begin{tabular}{lrrrrrr}
\toprule
System & Python & JavaScript & TypeScript & Go & Java & Rust \\
\midrule
Allfather + GPU reranker & \textbf{86.0} & \textbf{79.0} & 75.0 & 81.0 & 71.2 & \textbf{74.0} \\
Allfather, browser only & \textbf{86.0} & \textbf{79.0} & 73.0 & 80.0 & 66.2 & 66.0 \\
SweRankEmbed-Large & 81.0 & 52.0 & \textbf{81.0} & \textbf{82.0} & \textbf{73.8} & 68.0 \\
SweRankEmbed-Small & 79.0 & 52.0 & 75.0 & 69.0 & 58.8 & 60.0 \\
gpt-6-luna, one call & 51.0 & 37.0 & 33.0 & 56.0 & 35.0 & 26.0 \\
\midrule
Bugs & 100 & 100 & 100 & 100 & 80 & 100 \\
\bottomrule
\end{tabular}
\end{table}


\textbf{Language.} Table~\ref{tab:languages} breaks down hit@10 on all \NBugs{} bugs by language. \toolgpu is the most accurate system in JavaScript, Python, and Rust, and SweRankEmbed-Large in Go, Java, and TypeScript. The GPU reranker helps most in Rust and Java, where the browser ranker is weakest.

\textbf{What the report says.} The strata show where agents gain. When the report names a changed function (\res{\FF}{website-gpu}{sn-function} of the \NAgent{} bugs), \tool finds it \res{\FF}{website-gpu}{s-function}\% of the time and the gpt-6-luna agent \res{\FF}{codex-agent}{s-function}\%. When the report names neither a changed function nor a changed file (\res{\FF}{website-gpu}{sn-none} bugs), \tool reaches \res{\FF}{website-gpu}{s-none}\% and the agent \res{\FF}{codex-agent}{s-none}\%, while the Qwen3-Coder-Next agent drops to \res{\FF}{open-agent-qwen3-coder-next-fp8}{s-none}\% and Gemini~3.8 Flash asked once to \res{\FF}{gemini-single}{s-none}\%. The hand check of the strata gives Cohen's kappa \pending.

\begin{finding}{Answer to RQ4}
No system dominates across languages. Reports that name the faulty function are easy for every system; on reports that name nothing, single-call models and the open-weight agent fall far behind, while \tool and the free agent hold up.
\end{finding}

\subsection{The Reranker Across Languages}
\label{sec:reranker}

\begin{table}[t]
\centering
\small
\caption{The GPU reranker on training bugs, cross-validated with whole projects held out: function hit@10. Current pair: two adapters trained on Python only. A, B, C: one adapter trained on all six languages (A: the shipped recipe; B: larger groups, two epochs, language-balanced; C: B's recipe on a 4B model). $^*$ passes the rule (gain $\geq$ 2 points, Holm $p<0.05$).}
\label{tab:reranker}
\resizebox{\columnwidth}{!}{\begin{tabular}{lrrrrrr}
\toprule
Language & Bugs & No reranker & Python pair & A & B & C \\
\midrule
Python & 3174 & 79.3 & 81.9 & 83.6 & 82.7 & 82.9 \\
JavaScript & 438 & 68.3 & 72.8$^*$ & 72.8$^*$ & 72.8$^*$ & 74.4$^*$ \\
TypeScript & 586 & 56.3 & 59.2$^*$ & 61.3$^*$ & 60.1$^*$ & 58.7$^*$ \\
Go & 816 & 68.9 & 72.5$^*$ & 73.3$^*$ & 72.1$^*$ & 75.5$^*$ \\
Java & 768 & 62.9 & 68.0$^*$ & 70.0$^*$ & 68.6$^*$ & 58.1 \\
Rust & 120 & 73.3 & 80.8$^*$ & 83.3$^*$ & 84.2$^*$ & 85.0$^*$ \\
\bottomrule
\end{tabular}}
\end{table}


The GPU reranker was first trained on Python bugs only. Before scoring \fresh, I tested on training projects, cross-validated by whole projects, whether it transfers to the other five languages, and chose among three six-language candidates by a rule fixed in advance: a candidate passes a language if it improves hit@10 by at least two points with $p_\mathrm{Holm}<0.05$ (Table~\ref{tab:reranker}). The Python-trained reranker already passes in all five other languages. Candidate~A, a single adapter trained on all six languages with the original recipe, also passes in all five and is the one \toolgpu uses; candidate~C, a 4B-parameter model, collapsed on Java.

\subsection{Robustness}
\label{sec:robustness}

\begin{table}[t]
\centering
\small
\caption{Shortened (first 1,500 characters of the body) against full bug reports, function hit@10 on the same bugs.}
\label{tab:conditions}
\begin{adjustbox}{max width=\linewidth}\begin{tabular}{lrrrr}
\toprule
System & Bugs & Shortened & Full & Change \\
\midrule
gpt-6-astra, agent & 198 & 86.4 & 86.9 & +0.5 \\
\tool{}, browser only & 198 & 73.7 & 77.3 & +3.5 \\
gpt-6-luna, agent & 198 & 75.8 & 77.3 & +1.5 \\
\toolgpu{} & 198 & 74.2 & 76.8 & +2.5 \\
SweRankEmbed-Large & 198 & 70.2 & 72.7 & +2.5 \\
SweRankEmbed-Small & 198 & 62.6 & 62.1 & -0.5 \\
Qwen3-Coder-Next, agent & 198 & 60.6 & 61.6 & +1.0 \\
gpt-6-astra, one call & 198 & 56.1 & 60.6 & +4.5 \\
Gemini 3.8 Flash, one call & 198 & 51.5 & 55.6 & +4.0 \\
gpt-6-luna, one call & 198 & 37.9 & 42.9 & +5.1 \\
Qwen3-Coder-Next, one call & 198 & 33.3 & 40.9 & +7.6 \\
gpt-oss-120b, one call & 198 & 24.7 & 31.3 & +6.6 \\
\toolgpu{} & 580 & 76.0 & 77.9 & +1.9 \\
\tool{}, browser only & 580 & 71.6 & 75.3 & +3.8 \\
SweRankEmbed-Large & 580 & 70.9 & 72.9 & +2.1 \\
SweRankEmbed-Small & 580 & 65.3 & 65.9 & +0.5 \\
gpt-6-astra, one call & 580 & 50.0 & 55.0 & +5.0 \\
gpt-6-luna, one call & 580 & 32.8 & 39.8 & +7.1 \\
\bottomrule
\end{tabular}\end{adjustbox}
\end{table}


\textbf{Shortened reports.} Every system except SweRankEmbed-Small gains from the full report, by up to eight points (Table~\ref{tab:conditions}), and the order of the systems is essentially unchanged. Table~\ref{tab:short198} gives the full results on shortened reports, the condition on which the paid agents completed every bug.

\begin{table*}[t]
\centering
\caption{Shortened bug reports (title and first 1,500 characters of the body), the 198-bug agent subset, including the paid frontier agents. Columns as in Table~\ref{tab:full198}.}
\label{tab:short198}
\begin{tabular}{lcrrrrrrrr}
\toprule
System & Free & $n$ & hit@1 & hit@10 & 95\% CI & all@10 & MRR & Cost/bug & $\Delta$ hit@10 (W/L, $p_{\mathrm{Holm}}$) \\
\midrule
Claude Opus 5.5, agent &  & 198 & 77.8 & 90.9 & 86.2--95.1 & 49.5 & 0.828 & \$0.20 & +16.7 (41/8, $<$0.001) \\
gpt-6-astra, agent &  & 198 & 80.3 & 86.4 & 80.9--91.2 & 41.4 & 0.830 & \$0.29 & +12.1 (38/14, 0.038) \\
gpt-6-luna, agent & \checkmark & 198 & 67.2 & 75.8 & 70.0--81.2 & 34.3 & 0.710 & \$0.004 & +1.5 (26/23, 1.000) \\
\textbf{\toolgpu{}} & \checkmark & 198 & 45.5 & 74.2 & 68.3--79.9 & 43.4 & 0.540 & \$0.001 & reference \\
\tool{}, browser only & \checkmark & 198 & 41.9 & 73.7 & 68.4--79.4 & 39.4 & 0.520 & \$0 & -0.5 (7/8, 1.000) \\
SweRankEmbed-Large & \checkmark & 198 & 42.9 & 70.2 & 63.4--76.9 & 37.4 & 0.520 & -- & -4.0 (18/26, 1.000) \\
SweRankEmbed-Small & \checkmark & 198 & 39.4 & 62.6 & 54.6--70.2 & 32.8 & 0.456 & -- & -11.6 (16/39, 0.084) \\
Qwen3-Coder-Next, agent & \checkmark & 198 & 52.5 & 60.6 & 52.7--68.7 & 29.8 & 0.557 & \$0.028 & -13.6 (19/46, 0.036) \\
gpt-6-astra, one call &  & 198 & 43.4 & 56.1 & 47.8--64.4 & 21.7 & 0.489 & \$0.51 & -18.2 (13/49, $<$0.001) \\
Gemini 3.8 Flash, one call & \checkmark & 198 & 36.4 & 51.5 & 43.0--60.3 & 22.2 & 0.417 & \$0.11 & -22.7 (13/58, $<$0.001) \\
gpt-6-luna, one call & \checkmark & 198 & 30.3 & 37.9 & 29.6--46.0 & 15.7 & 0.338 & \$0.005 & -36.4 (4/76, $<$0.001) \\
Qwen3-Coder-Next, one call & \checkmark & 198 & 24.2 & 33.3 & 25.8--41.3 & 11.1 & 0.277 & \$0.007 & -40.9 (5/86, $<$0.001) \\
gpt-oss-120b, one call & \checkmark & 198 & 20.7 & 24.7 & 17.0--32.6 & 10.6 & 0.224 & \$0.006 & -49.5 (1/99, $<$0.001) \\
\bottomrule
\end{tabular}
\end{table*}


\textbf{Memorization.} On the SWE-bench Verified pilot, Claude Opus~5.5 asked once matched its own agent. On \fresh, gpt-6-astra asked once reaches \AstraSingleShort\% hit@10 on the bugs where its agent reaches \AstraAgentShort\%. A model that cannot read the code is far weaker on bugs it cannot have seen, which is what memorization of older benchmarks would predict. The pilot has only ten bugs, so this is suggestive rather than conclusive.

\textbf{Repeatability.} Asked again on 50 bugs, gpt-6-astra asked once changed its hit@10 outcome on \RepeatSingleChanged{} of \RepeatSingleN{} bugs, and its agent on \RepeatAgentChanged{} of \RepeatAgentN{}. One answer per bug is therefore a small source of noise compared with the differences reported above. \tool is deterministic.

\textbf{Long paths.} Some early agent runs on Windows silently dropped deeply nested files from their checkouts, never a file a fix changed. Repeating every comparison without the affected bugs (135 of \NBugs{}, 41 of \NAgent{}) changes no test's conclusion at $\alpha=0.05$. All later runs used Linux.

%% ---------------------------------------------------------------------------
\section{Discussion}
\label{sec:discussion}

\textbf{Implications for developers.} A free, local localizer is not a replacement for a frontier agent, but it is a strong first step. It finds a faulty function in its top ten as often as the best free agent, answers in seconds with the same answer every time, and never sends code off the machine. These properties suit uses where an agent is impractical: triaging every new issue in a large project, commenting on incoming issues automatically, or working on code that may not be sent to a model provider. For a developer fixing one bug who may share the code, a paid agent is more accurate and ranks the fault first more often.

\textbf{Implications for tool builders.} The escalation results suggest building the two as complements rather than competitors. A cheap localizer can answer the bugs it is confident about and call an agent for the rest; the gate in this study was fitted on development data and still selected the right bugs on unseen ones. Its ranking could also be handed to the agent as a starting point, which this study does not evaluate.

\textbf{Implications for researchers.} Evaluations on benchmarks that predate the models overstate models that may recall a fix. Comparing a model's single-call accuracy with its accuracy as an agent is a cheap diagnostic: on \fresh the two differ by about thirty points, on SWE-bench Verified they did not differ at all. Benchmarks should also report cost and the free baselines a developer actually has, because accuracy alone hides most of the trade-off.

%% ---------------------------------------------------------------------------
\section{Threats to Validity}
\label{sec:threats}

\textbf{Internal validity.} Closed models change without notice; I record each model's identifier and run date. The frontier agents include their vendors' system prompts, tools, and sub-agents, so I compare them as developers use them rather than as bare models. The open-weight agents run in my own harness; I fixed harness gaps found in run records before reading any answer, and a harness designed by their vendors might score higher. Each system answers each bug once; I measured the resulting noise on 50 bugs. Runs stopped by usage limits are reported on a random prefix.

\textbf{Construct validity.} The ground truth is what the developers' fix changed. A fix may include incidental edits or omit equally valid alternatives, and some issues describe missing features rather than defects. Hit@10 is a proxy for usefulness to a developer~\cite{kochhar2016practitioners,parnin2011automated}; I therefore also report hit@1 and all@10.

\textbf{External validity.} \fresh covers popular open-source projects in six languages over three months, and results may differ for proprietary code, other languages, or other kinds of bug. Every bug postdates every model's stated training cutoff, but older code in the same projects may be known to the models, which would favor them. \tool's own ranker was trained on projects disjoint from \fresh. AI-written issues or fixes that carry no marker could not be excluded.

\textbf{Conclusion validity.} The agent subset has \NAgent{} bugs, enough to detect a difference of about nine points; smaller differences are reported as ties. I correct for multiple comparisons within each table and cluster intervals by project.

%% ---------------------------------------------------------------------------
\section{Related Work}
\label{sec:related}

\textbf{Information-retrieval localization.} BugLocator combined a revised vector space model with similar fixed bugs~\cite{zhou2012buglocator}, BLUiR exploited the structure of code and reports~\cite{saha2013bluir}, and Ye et al. learned to rank files from many domain features~\cite{ye2014learning}. Bench4BL reproduced these techniques on 46 projects and showed that their performance varies widely across projects~\cite{lee2018bench4bl}. This line of work localizes files, mostly in Java; this study localizes functions in six languages.

\textbf{Neural localization.} Cross-encoder reranking~\cite{nogueira2019passage} and parameter-efficient fine-tuning~\cite{hu2022lora} underlie \toolgpu. SweRank trains a retriever and a reranker for issue localization and reports state-of-the-art results on SWE-bench Lite and LocBench~\cite{reddy2025swerank}; I compare against its published retrievers on fresh bugs.

\textbf{Agents and LLM-based localization.} SWE-agent introduced agent--computer interfaces for repositories~\cite{yang2024sweagent}, Agentless showed that a fixed hierarchical procedure is competitive with agents~\cite{xia2024agentless}, and LocAgent localizes with graph-guided search~\cite{chen2025locagent}. These systems are evaluated on SWE-bench and its variants, whose bugs predate current models. This study evaluates commercial agents as developers use them, on bugs that postdate them, with cost.

\textbf{Contamination and fresh benchmarks.} LiveCodeBench collects programming problems continuously after models' cutoffs to avoid contamination~\cite{jain2024livecodebench}. \fresh applies the same principle to bug localization and adds sealing and vetting for AI-written bugs.

\textbf{Developers and localization.} Kochhar et al. surveyed what practitioners expect from automated localization~\cite{kochhar2016practitioners}, and Parnin and Orso questioned whether ranked lists help programmers at all~\cite{parnin2011automated}. Both motivate reporting the top of the ranking, not only hit@10.

%% ---------------------------------------------------------------------------
\section{Conclusion}
\label{sec:conclusion}

On \NBugs{} bugs that no model could have seen, a free bug localizer that runs in the browser matches the best free coding agent at hit@10, outperforms open-weight models and every single-call model, and trails paid frontier agents, which a confidence gate can call only when needed. The gap between single-call and agent accuracy suggests that older benchmarks overstate what models can do on new bugs. \fresh, every run record, and the protocol with its amendments are available for others to extend with new systems and new bugs.

\section*{Data Availability}

The benchmark (bugs, ground truth, and both report conditions), every system's answers and run records, the prompts and harnesses, the analysis scripts, and the protocol with its dated amendments are available at \artifact. Every table, figure, and number in this paper is generated from the committed run records by a single script.

\begin{acks}
This research was supported in part by research cyberinfrastructure resources and services provided by the Partnership for an Advanced Computing Environment (PACE) at the Georgia Institute of Technology, Atlanta, Georgia, USA. AI disclosure: \pending{} (to be written according to the ACM Policy on Authorship).
\end{acks}

\bibliographystyle{ACM-Reference-Format}
\bibliography{refs}

\end{document}
